\documentclass{article}


% if you need to pass options to natbib, use, e.g.:
%     \PassOptionsToPackage{numbers, compress}{natbib}
% before loading neurips_2025


% ready for submission
\usepackage{neurips_2025}

% to compile a preprint version, e.g., for submission to arXiv, add add the
% [preprint] option:
%     \usepackage[preprint]{neurips_2025}


% to compile a camera-ready version, add the [final] option, e.g.:
%     \usepackage[final]{neurips_2025}


% to avoid loading the natbib package, add option nonatbib:
%    \usepackage[nonatbib]{neurips_2025}


\usepackage[utf8]{inputenc} % allow utf-8 input
\usepackage[T1]{fontenc}    % use 8-bit T1 fonts
\usepackage{hyperref}       % hyperlinks
\usepackage{url}            % simple URL typesetting
\usepackage{booktabs}       % professional-quality tables
\usepackage{amsfonts}       % blackboard math symbols
\usepackage{nicefrac}       % compact symbols for 1/2, etc.
\usepackage{microtype}      % microtypography
\usepackage{xcolor}         % colors
\usepackage{amsmath}
\usepackage{bm}
\usepackage{wrapfig}
\usepackage{comment}
\usepackage[noend,linesnumbered,ruled,vlined]{algorithm2e}
\usepackage{graphicx}
\usepackage{amssymb}
\usepackage{wrapfig}
\usepackage{multirow}
\usepackage{comment}
\title{Foveated Segment Anything}
\usepackage{caption}
\usepackage{tabularx,colortbl,xcolor}

% The \author macro works with any number of authors. There are two commands
% used to separate the names and addresses of multiple authors: \And and \AND.
%
% Using \And between authors leaves it to LaTeX to determine where to break the
% lines. Using \AND forces a line break at that point. So, if LaTeX puts 3 of 4
% authors names on the first line, and the last on the second line, try using
% \AND instead of \And before the third author name.


\author{%
  David S.~Hippocampus\thanks{Use footnote for providing further information
    about author (webpage, alternative address)---\emph{not} for acknowledging
    funding agencies.} \\
  Department of Computer Science\\
  Cranberry-Lemon University\\
  Pittsburgh, PA 15213 \\
  \texttt{hippo@cs.cranberry-lemon.edu} \\
  % examples of more authors
  % \And
  % Coauthor \\
  % Affiliation \\
  % Address \\
  % \texttt{email} \\
  % \AND
  % Coauthor \\
  % Affiliation \\
  % Address \\
  % \texttt{email} \\
  % \And
  % Coauthor \\
  % Affiliation \\
  % Address \\
  % \texttt{email} \\
  % \And
  % Coauthor \\
  % Affiliation \\
  % Address \\
  % \texttt{email} \\
}


\begin{document}


\maketitle


\begin{abstract}
Augmented Reality (AR) encompasses transformative technologies that are redefining how humans interact with their environment. A key component of AR is image segmentation, which breaks down the user's front-view scene into distinct regions for analysis. This process is essential for accurately overlaying digital content onto the physical world by detecting and isolating relevant objects. However, despite its importance, image segmentation poses significant computational demands and latency issues on AR devices, which can severely impact the overall user experience. In this paper, we propose~\textit{Foveated Segment Anything Model} (FoSAM), a framework built upon the Segment Anything Model (SAM) that utilizes real-time gaze data to focus segmentation on regions of interest, substantially lowering computational cost. Experimental results show that FoSAM reduces computational cost by over $50\times$, enabling a seamless visual experience for users, as confirmed by our real-world user study. %(\textcolor{red}{https://github.com/ElecBeholder/FoSAM})

\end{abstract}


\section{Introduction}
\begin{wrapfigure}{r}{0.5\textwidth}
  \begin{center}
    \includegraphics[width=0.5\textwidth]{figures/foveated_example.pdf}
  \end{center}
  \caption{Application of VLM in AR.}
  \label{fig:ar-application}
\end{wrapfigure}
Image segmentation~\cite{long2015fully, badrinarayanan2017segnet, kirillov2023segment, xie2021segformer} is a core task in computer vision that involves dividing an image into meaningful regions to support visual content analysis and interpretation. Building on this, instance segmentation~\cite{he2017mask, bolya2019yolact, wang2020solov2, yang2019video} identifies and outlines each individual object within a scene. This capability is particularly crucial in augmented reality, where accurate object detection and separation enable precise interaction, manipulation, and seamless overlay of virtual content onto the physical environment, enhancing both immersion and contextual awareness. 

Instance segmentation serves as a foundational component for numerous AR applications. For example, in an AR-assisted grocery shopping scenario, real-time segmentation enables users to identify products on a shelf as they look at them (Figure~\ref{fig:ar-application} (a)). In educational contexts (Figure~\ref{fig:ar-application} (b)), segmentation can detect individual components of complex diagrams the user is viewing, enhancing student engagement and understanding. These segmented objects can also be passed to downstream applications (e.g., vision-language models (VLMs)), to provide detailed explanations or context-aware information. Furthermore, segmentation enables direct object manipulation, allowing users to edit or interact with specific elements in their environment, creating a seamless bridge between the physical and virtual worlds (Figure~\ref{fig:ar-application} (c)). Additional AR use cases leveraging image segmentation are discussed in~\cite{gonzalez2019applications, gonzalez2020nextmed, tanzi2021real, alhaija2017augmented}.

The Segment Anything Model (SAM)~\cite{kirillov2023segment, ravi2024sam} is among the most advanced models for image segmentation today. However, despite its strong performance, SAM is computationally intensive, making it difficult to deploy on resource-limited AR devices that handle high-resolution imagery. This results in high overhead and latency, degrading overall system performance and user experience.

Unlike typical image processing scenarios, AR device users display distinct behavioral patterns. They tend to focus on small, specific areas within a scene before shifting their attention elsewhere. For instance, as shown in Figure~\ref{fig:eye_study} (a), a user wearing AR glasses may first focus on the door (left) and then turn their head to look at the bookshelf (right). This behavior naturally segments the video stream into distinct video fragments (VF) based on head movements. In the first VF, where the user’s gaze remains on the door across several frames, the segmentation results can be reused, avoiding redundant computation. Similarly, in the second VF, segmentation efforts can be focused solely on the bookshelf.
Moreover, as illustrated in Figure~\ref{fig:ar-application}, it is often beneficial to restrict segmentation to the instances of interest (IOI) currently under the user’s gaze, as this typically reflects the user’s active focus. This leads to a more efficient approach for AR image segmentation: \textbf{prioritize processing for gaze-identified IOIs while ignoring non-essential regions}. This strategy aligns with the principles of \textit{foveated rendering}~\cite{patney2016towards}, which improves rendering efficiency by showing full-resolution detail only within the user’s line of sight, reducing visual fidelity in peripheral areas to save computations.

In this work, we aim to reduce the high computational cost of segmentation in AR by leveraging natural human eye movement. We introduce a novel segmentation framework based on Efficient SAM~\cite{xiong2024efficientsam} (ESAM), a lightweight version of the original SAM that maintains nearly the same performance. By limiting segmentation to only the regions the user is actively viewing, our proposed variant, called \textit{foveated SAM} (FoSAM), achieves more than a $50\times$ reduction in processing latency without compromising the user's visual experience. Our contribution can be summarized as follows:
\begin{itemize}
\item We introduce a novel and important perspective on instance segmentation with SAM that exploits human eye behavior to reduce computational costs in AR settings.
\item We propose~\textit{FoSAM}, a lightweight segmentation framework built on ESAM, which processes high-resolution input images and performs instance segmentation on the IOI with extremely low computational cost.
\item Building on FoSAM, we introduce~\textit{FoSAM Streaming Algorithm} (FSA), an efficient instance segmentation framework tailored for real-time AR/VR applications. FSA, exploits temporal continuity across frames and human gaze patterns to optimize segmentation, delivering enhanced performance in dynamic AR environments.

\end{itemize}
\begin{figure}
    \centering
    \includegraphics[width=1\columnwidth]{figures/eye_study_new.pdf}
    \caption{(a) An example on gaze movement for the AR scenario. (b)
Trace of eye gaze within a segment. (c) Meta Orion AR glass. (d) Meta Quest Pro VR headset.}
    \label{fig:eye_study}
\end{figure}

\section{Background and Related Work}
\label{sec:related_work}
\subsection{Human Eye Behavior in AR/VR Environment}
\label{sec:bg:eye-behavior}
The human eye operates in three main modes of movement, each serving distinct functions: fixation, where the eye remains stationary and focuses on a single point; saccadic movements, rapid, jerky shifts in gaze from one target to another; and smooth pursuit, where the eye follows a moving object in a smooth manner. Among these, smooth pursuit is less frequent, while fixation and saccadic movements dominate the majority of visual activity, as illustrated in Figure~\ref{fig:eye_study} (b).
During fixation, the gaze remains locked on a single point, and visual acuity is concentrated around this focal region, which is densely populated with photoreceptor cells. This region, although small, has the sharpest vision, with acuity rapidly diminishing as we move away from the gaze point, resulting in less sensitivity to fine details in peripheral vision.
In addition to fixation, the eye engages in saccadic motion, quickly jumping from one point of focus to another. These rapid movements occur in fractions of a second, enabling the brain to rapidly gather visual information from different areas of the field. 
%During a saccade, the visual system's sensitivity is temporarily reduced, a phenomenon known as~\textit{saccadic suppression}~\cite{matin1974saccadic, krekelberg2010saccadic}. This decrease in sensitivity helps prevent the brain from perceiving the blur caused by the swift movement of the eyes.

%\subsection{Gaze Behavior in AR/VR Environment}
%\label{sec:bg:gaze-behavior-arvr}


\begin{wrapfigure}{r}{0.5\textwidth}
    \centering
    \includegraphics[width=0.5\columnwidth]{figures/trace_new.pdf}
    \caption{(a) Video Fragments in the user's view. (b) Distances (in pixels) between consecutive gaze locations. (c) Distribution of the gaze distances. (d) Pixel differences of consecutive frames. (e) Distribution of the pixel differences.}
    \label{fig:aria_study}
\end{wrapfigure}

Figure~\ref{fig:eye_study} (c) and (d) present the Meta Aria AR glasses~\cite{engel2023project} and Meta Quest Pro~\cite{quest_pro}, respectively. Both devices feature an outward-facing camera that continuously captures high-resolution images of the user's front view, along with an inward-facing eye-tracking (ET) camera that records monochrome images of the user's eyes. Most modern AR devices are equipped with eye trackers that takes the eye images and provide highly precise gaze direction estimates with minimal latency.
To investigate human gaze behavior while using AR devices, we perform an in-depth analysis using the Aria Everyday Activities Dataset~\cite{lv2024ariaeverydayactivitiesdataset}. This dataset contains 143 sequences of frames captured by the AR device, which align with the user's field of view and track gaze movements across each frame in different environments, as shown in Figure~\ref{fig:aria_study}.

As the user adjusts their head orientation, the front view also changes. To quantify these head movements, we calculate the image difference by measuring the Euclidean distance between corresponding pixels of consecutive frames. Small pixel differences between frames indicate that the user is maintaining a stable head orientation while observing a static scene. If the pixel difference is below a certain threshold, the frames are considered highly similar and nearly indistinguishable to the human eye. These frames are grouped into a single video fragment (VF), as shown in Figure~\ref{fig:eye_study} (a) and Figure~\ref{fig:aria_study} (a). Figure~\ref{fig:aria_study} (d) presents the percentage of pixel differences over a sample time interval spanning four VF, with consecutive frames grouped based on a threshold of pixel changes, marked by the yellow line. As shown in Figure~\ref{fig:aria_study} (e), $32\%$ of consecutive frames exhibit less than $5\%$ pixel value changes,~\textbf{suggesting a high temporal similarity between the consecutive frames and potential for reusing segmentation results across frames}.

Furthermore, within each VF, segmentation results can be reused if the gaze remains relatively stable and consistently points to the same IOI. To support this, we analyze the distances (in pixels) between consecutive gaze locations within a VF, as shown in Figure~\ref{fig:aria_study} (b). Our analysis reveals that a threshold of \(20\) pixels effectively groups gaze locations during the fixation phase, where the user focuses on a single IOI. Gaze distances exceeding this threshold indicate a saccade, as seen with the rapid gaze changes in VF3 of Figure~\ref{fig:aria_study} (a). As shown in Figure~\ref{fig:aria_study} (c), \(87\%\) of the frames within each VF have a gaze distance of less than \(20\) pixels, indicating that AR users typically focus on one or two IOIs during each VF. This provides an opportunity to enhance image segmentation efficiency by~\textbf{focusing processing on the IOI and reusing segmentation results when gaze shifts are minimal}.


\subsection{Literature Review on Segmentation}
\label{sec:bg:segmentation}  
Semantic segmentation~\cite{Shelhamer2014FullyCN, Badrinarayanan2015SegNetAD, chen2017deeplab, Touvron2020TrainingDI, Cheng2021PerPixelCI, Minaee2020ImageSU, Li2023TransformerBasedVS, Xiong2019UPSNetAU, Cheng2019PanopticDeepLabAS} is a key task in computer vision that involves partitioning an image into distinct regions or segments to simplify content analysis and interpretation. A more advanced form of this task is instance segmentation~\cite{He2017MaskR, Li2020FullyCN, Neven2019InstanceSB, Brabandere2017SemanticIS}, which aims to distinguish individual instances of the same object class. Unlike semantic segmentation, which labels each pixel without differentiating between instances of the same object class, instance segmentation provides a finer level of detail by distinguishing each instance.
To improve segmentation efficiency, previous research has focused on developing learnable input downsampling techniques that adjust sampling resolution in a selective manner. In \cite{recasens2018learning}, the authors propose a saliency-based distortion layer for convolutional neural networks that enhances spatial sampling of input data in image classification tasks. Subsequent works, such as \cite{jin2021learning, thavamani2021fovea, marin2019efficient,zeng2025foveated}, follow similar approaches by learning a saliency score for each pixel to guide the downsampling process.
%leading to improved performance in semantic segmentation tasks. 
%LZU~\cite{thavamani2023learning} offers an efficient method to learn both zoom and unzoom processes. 
Previous approaches perform image downsampling at the pixel level using Convolutional Neural Networks (CNNs). In contrast, FoSAM is a transformer-based model that adaptively downsamples the input image at the token level, guided by the user's gaze location and the shape of the IOI.
Furthermore, previous studies~\cite{yang2019video, wang2021end, yan2023universal, rajivc2023segment, lin2021video, wu2022seqformer} on video instance segmentation process consecutive frames together.
%(e.g., VisTR~\cite{yan2023universal} and SeqFormer~\cite{wu2022seqformer} process five consecutive frames at once). 
This methods leverage temporal correlations across frames to improve performance; however, it introduces considerable latency, as processing can only begin once all frames are available.

\subsection{Segmentation Task Latency}
\begin{comment}
\begin{wraptable}{r}{0.3\textwidth}
    \centering\resizebox{0.3\columnwidth}{!}{
       \begin{tabular}{c|c|c}
        \toprule
    Model  & Image Size & Latency \\ \midrule
    SAM-B & $1024^2$  & 1897.4 ms\\ \midrule
    \multirow{2}{*}{ESAM-S} & $1024^2$  & 600.3 ms\\
     & $640^2$  &  102.6ms\\ \midrule
    \multirow{2}{*}{ESAM-T} & $1024^2$  & 279.7 ms\\    
     & $640^2$  & 46.7 ms\\ 
    \bottomrule
    \end{tabular}
}
\caption{Processing latencies of segmentation models on an edge GPU. (\textcolor{red}{change it to 640x640 and smaller})}
\label{tab:reduce_size}
% \vspace*{-1\baselineskip}
\end{wraptable}
\end{comment}

\begin{wraptable}{r}{0.4\textwidth}
    \centering\resizebox{0.4\columnwidth}{!}{
       \begin{tabular}{c|c|c|c}
        \toprule
    Image Size & SAM-B  & ESAM-S & ESAM-T \\ \midrule
    $1024^2$ & 1897.4 ms & 600.3 ms  & 279.7 ms\\
    $640^2$ & 441.8 ms & 102.6 ms  & 46.7 ms\\
    \bottomrule
    \end{tabular}
}

\caption{Processing latencies of segmentation models on AR platform.}
\label{tab:reduce_size}
% \vspace*{-1\baselineskip}
\end{wraptable}

Despite its significance, the segmentation task presents considerable computational challenges, especially on resource-constrained AR devices, primarily because of the high resolution of images these devices capture. This heavy data load leads to significant computational latency, severely limiting performance and responsiveness.
To explore this issue, we evaluate the processing latency of the image segmentation task using well-established neural networks, including Segment Anything Base model (SAM-B)~\cite{kirillov2023segment, ravi2024sam}, Efficient SAM Tiny (ESAM-T)~\cite{xiong2024efficientsam}, and ESAM Small (ESAM-S). The latencies are profiled on the Qualcomm XR2 Gen2 platform~\cite{qualcomm_xr2_gen2} using the Qualcomm AI Hub toolkit~\cite{qualcomm_aihub}, which is equipped by Meta Quest Pro~\cite{quest_pro} and Meta Aria AR glasses~\cite{engel2023project}.
Table~\ref{tab:reduce_size} shows that the processing latencies for $1024 \times 1024$ input images are 1897.4ms, 600.3ms, and 279.7ms for SAM-base (SAM-B), Efficient SAM-small (ESAM-S), and Efficient SAM-tiny (ESAM-T), respectively. Furthermore, even with a smaller input resolution of $640 \times 640$, the latency still ranges from 441.8ms to 46.7ms. These latency values fall short of the threshold required for a seamless visual experience, as prior studies indicate that latencies below 30 ms are essential for maintaining optimal visual fluidity~\cite{kaaresoja2014towards, visual_requirement,albert2017latency}. 

\section{Methodology}
\label{sec:Methodology}
\begin{wrapfigure}{r}{0.5\textwidth}
  \begin{center}
    \includegraphics[width=0.5\textwidth]{figures/fosam_overall.pdf}
  \end{center}
  \caption{System Deployment of FoSAM.}
  \label{fig:fosam_framework}
\end{wrapfigure}
Figure~\ref{fig:fosam_framework} illustrates the computational pipeline of the FoSAM framework. During runtime, the inward-facing sensor of the AR/VR device continuously captures images of the user’s eye and sends them to the gaze tracker, which estimates the gaze direction instantaneously with high accuracy (5 to 10 milliseconds)~\cite{stein2021comparison, hou2024unveiling}. At timestep t, this estimated gaze direction $g_{t}$, along with a high-resolution image $I_{t}$ captured by the front-facing camera, is provided as input to FoSAM. FoSAM then produces a segmentation map $M_{t}$ focused on the IOI, which can be reused across frames where the gaze location remains similar. To simplify notation, we omit the subscript $t$ in the following sections.

The core idea of FoSAM is to employ a lightweight saliency encoder to identify and select a subset of image tokens, which are then passed to the segmentation network. By processing only the selected tokens, the segmentation network significantly reduces its computational overhead.

%Next, we begin by reviewing the Efficient SAM (ESAM)~\cite{xiong2024efficientsam} baseline to establish the foundation for our design (Section~\ref{sec:Preliminaries}). Next, we describe the saliency prediction module and the dynamic window adjustment mechanism that enable efficient and adaptive processing (Section~\ref{sec:saliency-design} and Section~\ref{sec:dynamic-window}). Finally, we detail the training methodology, including the design of the loss functions used to optimize the model (Section~\ref{sec:training}). For clarity, we omit the timestep notation $t$ as the discussion pertains to operations performed within each individual timestep.

\subsection{Preliminaries}
\label{sec:Preliminaries}

FoSAM builds on ESAM~\cite{xiong2024efficientsam}, which includes three main components: an~\textit{image encoder} $E(.)$, a~\textit{prompt encoder} $R(.)$, and a~\textit{mask decoder} $D(.)$. Given an input RGB image $I \in \mathbb{R}^{H \times W \times 3}$, it is first tokenized and positionally embedded into a set of tokens $X = \{x_{ij}\}$ with a size of $H_e \times W_e$, where each token $x_{ij} \in \mathbb{R}^{1 \times d_t}$. $H_e$ and $W_{e}$ equal $\frac{H}{P}$ and $\frac{W}{P}$, where P is the patch size of the tokenizer. These tokens are passed through the image encoder to produce $Y=\{y_{ij}\}=E(X)$, which retains the same spatial resolution. In parallel, a point-based gaze prompt $g$ is processed by the prompt encoder $R(.)$, and its output is fused with $Y$ in the decoder to generate the final segmentation mask $M \in \mathbb{R}^{H \times W}$. In ESAM, the image encoder $E(.)$ is a 12-layer Vision Transformer, and the decoder $D(.)$ consists of a pre-trained two-layer CNN and a two-layer ViT. Since $E(.)$ dominates the computational cost, we focus on optimizing it by introducing a~\textit{gaze-guided saliency encoder} $S(.)$, which selects a subset of $K$ important tokens from $X$ based on the gaze direction $g$, reducing the input size to $E(.)$ from $H_e \times W_e$ to $K$ and greatly decrease the computational cost of the segmentation task. The entire selection process is end-to-end trainable, resulting in strong accuracy performance. Next, we describe the design of gaze-guided saliency encoder $S(.)$.



%FoSAM is based on ESAM~\cite{xiong2024efficientsam}, which consists of three modules: an~\textit{image encoder}, a~\textit{prompt encoder} and a~\textit{mask decoder}, denoted as \(f^{\text{ESAM}}\), \(h^{\text{ESAM}}\) and \(g^{\text{ESAM}}\), respectively. Given an input image \(\bm{I} \in \mathbb{R}^{H^I \times W^I}\), we first perform tokenisation and positional embedding to obtain \(H^t \times W^t\) tokens \(\{\bm{x}_i\}_{i=1}^{H^t \times W^t} \in \mathbb{R}^{d^t},\) where \(d^t\) denotes the token dimension. The image encoder then encodes all tokens. The encoded tokens are denoted as \(\{\bm{y}_i\}_{i=1}^{H^t \times W^t} \in \mathbb{R}^{d^t}=f^{\text{ESAM}}(\{\bm{x}_i\}),\) which are subsequently fed into the mask decoder. Simultaneously, the gaze point, treated as a point-type prompt, is passed through the prompt encoder, and the two streams are fused to predict the final segmentation mask. In ESAM, the image encoder is composed of a 12-layer Vision Transformer (ViT), while the mask decoder includes a pre-trained two-layer CNN followed by a two-layer ViT. As the image encoder accounts for the majority of the computational load in ESAM, it becomes the primary focus of our optimization efforts. Specifically, we propose a~\textit{gaze-guided saliency encoder} designed to reduce the number of input tokens to the image encoder from $H^t \times W^t$ to a significantly smaller target size $K$, where $K \ll H^t \times W^t$.

\subsection{FoSAM framework}
\label{sec:fosam_frameowrk}
\subsubsection{Gaze-Guided Saliency Encoder}
\label{sec:saliency-design}
\begin{figure}
    \centering
    \includegraphics[width=1\columnwidth]{figures/fosam_fig.pdf}
    \caption{An overview of FoSAM framework.}
    \label{fig:fosam_fig}
\end{figure}
The gaze-guided saliency encoder is designed to assign a saliency score to each token $X$ based on the user's gaze direction $g$. A common approach predicts token-wise saliency scores independently for each token~\cite{liu2021visual,liu2023tracking}, and then removes those with low scores. However, this strategy often leads to poor accuracy because it lacks spatial coherence. The independently predicted scores can result in fragmented saliency maps with missing regions, which negatively affect the segmentation quality of the image decoder. Additionally, this approach requires computing a saliency score for every token, resulting in computational complexity that grows linearly with the number of tokens. This significantly limits the potential for reducing model size and improving efficiency.

To address these limitations, we propose estimating the saliency scores of tokens using a 2D Gaussian distribution. The top portion of Figure~\ref{fig:fosam_fig} illustrates the overall pipeline, with the black and blue arrows show the inference flow. The input image $I$ is first uniformly downsampled from $H \times W$ to $\frac{H}{v} \times \frac{W}{v}$, resulting in a lower-resolution image $I'$, where $v$ denotes the downsampling ratio. This step effectively reduces the computational cost of the saliency function $S(.)$. Next, $I'$ is passed through a pre-trained ESAM tokenizer and a positional embedding module to generate token features $X' = \{x'_{ij}\}$ with dimensions $H'_e \times W'_e$. $X'$ is then processed by the saliency module $S(.)$, which outputs $Z = S(X')$ containing the same number of tokens. Tokens $Z$ are selected based on their spatial proximity to the gaze point $g$ obtained from the gaze tracker. The token nearest to $g$, denoted as $z_g \in Z$, captures essential information about the IOI, such as its shape and scale, and serves as the foundation for predicting a bounding window that encloses the IOI. 

To simplify the window generation process while preserving spatial continuity, we introduce a Gaussian predictor to estimate the saliency score of each token $x\in X$. This predictor, $G(\cdot)$, takes $z_g$ as input and outputs the parameters of a 2D Gaussian distribution: the mean $(\mu_x, \mu_y)$, standard deviations $(\sigma_x, \sigma_y)$, and correlation coefficient $\rho$. Let $f_{G(z_{g})}(.,.)$ represent the probability density function of Gaussian distribution with the parameter specified by $G(z_{g})$. The saliency score for each token $x \in X$ is then computed as its likelihood under the predicted 2D Gaussian distribution, denoted as $f_{G(z_{g})}(c_{i}(x),c_{j}(x))$, where $c_{i}(x)$ and $c_{j}(x)$ is the 2D coordinate of $x$ within the token map.


%To address this objective, it is essential to retain tokens that are critical for the segmentation task while discarding those that are less relevant. A common strategy involves learning token-wise saliency scores [cite TODO], which quantify the importance of each token for the downstream task, enabling the removal of low-scoring tokens. However, this approach faces several limitations in a gaze-guided pruning context. First, due to architectural constraints, the model cannot directly leverage gaze point coordinates and instead requires an additional encoding module, similar to the prompt encoder used in ESAM. Second, because token scores are predicted independently, spatial continuity is not guaranteed, often resulting in fragmented token maps with “holes” that degrade segmentation quality. Third, the saliency encoder must compute a score for every token, leading to computational complexity that scales linearly with token count, hindering efforts to reduce model size and improve efficiency.

%To address these limitations, we propose using a 2D Gaussian distribution to estimate the saliency scores of the tokens associated with the IOI. Figure~\ref{fig:fosam_fig} (top part) details this pipeline. First, the input image is evenly downsampled from \(\bm{I} \in \mathbb{R}^{H^{I} \times W^{I}}\) to \(\tilde{\bm{I}} \in \mathbb{R}^{\tilde{H}^{I} \times \tilde{W}^{I}}\), this step can substantially reduce the cost of the saliency encoder. Next, the two images are passed separately through the pre‑trained ESAM tokenizer and positional embedding module, producing two sets of tokens: \(\{\bm{x}_{i}\}_{i=1}^{H^{t}\!\times W^{t}} \in \mathbb{R}^{d^t}\) for the original, and \(\{\bm{\tilde{x}}_{i}\}_{i=1}^{\tilde{H}^{t} \times \tilde{W}^{t}} \in \mathbb{R}^{d^t}\) for the downsampled. Then, \(\{\bm{\tilde{x}}_i\}\) is passed to our lightweighted saliency encoder \(f^{\text{SA}}\), which produces token embeddings \(\{\bm{\tilde{z}}_i\}=f^{\text{SA}}(\{\bm{\tilde{x}}_{i}\})\).

%The output token embeddings $\{\bm{\tilde{z}}_i\}$ are selected based on their spatial proximity to the gaze location provided by the gaze tracker. The token closest to the gaze point is denoted as $\bm{\tilde{z}}_g \in \{\bm{\tilde{z}}_i\}$, which captures key information, such as shape and scale, about the IOI, and serves as the basis for predicting a window that encloses it. To simplify the window generation process and maintain spatial continuity, we introduce a Gaussian predictor that estimates a saliency score for each token. This predictor takes $\bm{\tilde{z}}_g$ as input and outputs the parameters of a 2D Gaussian distribution $G(\cdot)$, including the mean $(\mu_x, \mu_y)$, standard deviations $(\sigma_x, \sigma_y)$, and correlation coefficient $\rho$. The saliency score for each token $x_{i}$ is then derived from its probability under the predicted 2D Gaussian distribution, denoted as $G(i)$.


\subsubsection{Adaptive Token Sampling Mechanism}
\label{sec:dynamic-window}
Based on the saliency scores generated by the 2D Gaussian map, only a subset of tokens with the highest scores are retained. Specifically, a larger IOI results in more tokens receiving high saliency scores, corresponding to a Gaussian distribution with lower variance. Conversely, smaller IOIs produce higher-variance distributions and fewer high-saliency tokens. As a result, we retain the top $K$ tokens, where $K$ is determined adaptively based on the Gaussian distribution. Specifically, let $\tau \in [0,1]$ denotes a predefined threshold of the token elimination, the token $x_{i}$ are kept only if specified as follows:
\begin{equation}
\label{eqn:thres}
\small
f_{G(z_{g})}(c_{i}(x), c_{j}(x)) \geq \max_{x\in X}[f_{G(z_{g})}(c_{i}(x), c_{j}(x))]\times \tau  
\end{equation}
Here, $\tau$ denotes the threshold applied to the saliency scores. Additionally, we define a minimum token count $\epsilon$ to ensure a sufficient number of tokens are selected. If the number of tokens satisfying the condition in Equation~\ref{eqn:thres} falls below $\epsilon$, the top-$\epsilon$ tokens are selected based on their saliency scores. By setting $\tau$ to a predefined value (e.g., 0.01), Equation~\ref{eqn:thres} enables adaptive token selection, allowing the number of retained tokens to vary according to the size of the IOI. Let $X_{s}\subseteq X$ represent set of tokens selected under equation~\ref{eqn:thres}.
The selected tokens are passed through the image encoder, yielding $E(X_s)$, which is then fed into the mask decoder $D(.)$ as shown by the low portion of Figure~\ref{fig:fosam_fig}. As noted in Section~\ref{sec:Preliminaries}, $D(.)$ expects a fixed input size of $H_e \times W_e$. To satisfy this requirement, $E(X_s)$ is zero-padded to the target dimensions, the results is denoted as $Y_s$. The decoder then generates the predicted segmentation mask $M_{pred}$. The following section details the training procedure.

\begin{comment}
After getting the saliency score for each token, a typical pruning strategy is to retain a fixed number of tokens via a top‑\(K\) selection. In practice, however, object size varies widely, so a fixed window size is inefficient; instead, the window should adapt to the object scale. In this work, we use a 2-D Gaussian distribution as the pruning window. The object size is intuitively reflected in the Gaussian standard deviation~\(\sigma\), and consequently in the density function~\(G\). Therefore, we define our selected tokens \(\{\bm{x}_i\}^s\) that are retained by dynamic window pruning as

\begin{equation}
\begin{aligned}
&\tau = \max(G)\times t, \\
&\{\bm{x}_i\}^s =
\begin{cases}
  \bigl\{\bm{x}_i \,\bigl|\, G(h_i, w_i) > \tau \bigr\}, &
  \text{if } \bigl|\,\{\bm{x}_i\}^{s}\bigr| > \varepsilon, \\[6pt]
  \operatorname{top}\varepsilon~\bigl(\{\bm{x_i}\};\, G\bigr), &
  \text{otherwise}.
\end{cases},\\
&\mathbb{E}\Bigl[\bigl\lvert \{\bm{x}_i\}^{s} \bigr\rvert\Bigr] \approx K,
\end{aligned}
\end{equation}

where \(\tau\) is the dynamic saliency threshold and \(t\in [0,1]\) is a fixed cut ratio. When \(t\) increased, the overall window size \(\bigl|\,\{\bm{x}_i\}^{s}\bigr|\) shrinks accordingly. The \(\operatorname{top}\varepsilon~\bigl(\{\bm{x_i}\};\, G\bigr)\) denotes the set of \(\varepsilon\) tokens with largest \(G(h_i,w_i)\), and \(\varepsilon\) is the minimum window size. Since large/small IOI tend to have larger/smaller~\(\sigma\) and thereby larger/smaller \(\max(G)\) (as dictated by the 2‑D Gaussian possibility density formula,) our pruning mechanism can adapt dynamically to the target object's scale, while the cut ratio~\(t\) governs the overall window size.

%After getting \(\{\bm{x}_i\}^{s}\), we encode them with \(f^{\text{ESAM}}\) to obtain \(\{\bm{y}_i\}^{s}\). However, as noted in Section~\ref{sec:Preliminaries}, the subsequent module \(g^{\text{ESAM}}\) contains convolutional layers and therefore requires a fixed input size of \(H^{t}\times W^{t}\) tokens. So, we zero‑pad \(\{\bm{y}_i\}^{s}\) to \(H^{t}\times W^{t}\) while keeping real tokens \(\{\bm{y}_i\}^{s}\) at their original positions. Finally, we can get the predicted segmentation mask from \(g^{\text{ESAM}}\), and next part we will introduce the detailed training procedure.
\end{comment}


\subsubsection{Training and Loss Function}
\label{sec:training}

Figure~\ref{fig:fosam_fig} illustrates the training and inference workflows of FoSAM, with data flows for inference and training shown in black+blue and black+grey, respectively. During inference, only the token $z_g$, which is closest to the gaze location $g$, is used to predict the saliency score. Let $M_{gt}$ represent the ground-truth binary mask corresponding to the input $I$, where $M_{gt, ij} = 1$ indicates that the pixel lies within the IOI, and 0 otherwise. Additionally, let $X_{M_{gt}} \subseteq X$ denote the subset of tokens whose spatial positions fall within the IOI as defined by $M_{gt}$, and let $Z_{M_{gt}} = S(X_{M_{gt}})$ represent the outputs of the $S(.)$ of the corresponding tokens.
The loss function used during training comprises three components: the negative log-likelihood, \(\mathcal{L}_{\mathrm{nll}}\), the classification loss \(\mathcal{L}_{\mathrm{cls}}\), and the Dice loss \(\mathcal{L}_{\mathrm{dice}}\). The objective of $\mathcal{L}_{\mathrm{nll}}$ is to train the saliency module $S(\cdot)$ and the Gaussian predictor $G(\cdot)$ to generate a 2D Gaussian map $f_{G(S(x))}(\cdot)$, where $x \in X_{M_{gt}}$, that closely approximates the ground-truth mask $M_{gt}$, as detailed below:
\begin{equation}
\label{eqn:loss1}
\small
\mathcal{L}_{\mathrm{nll}} = -\sum_{x_{m}\in X_{M_{gt}}} \sum_{x'_{m}\in X_{M_{gt}}} log [f_{G(S(x_{m}))}(c_{i}(x'_{m}), c_{j}(x'_{m}))]
\end{equation}
where $\sum_{x'_{m}\in X_{M_{gt}}} log [f_{G(S(x_{m}))}(c_{i}(x'_{m}), c_{j}(x'_{m}))]$ represents the sum of the log-probabilities for each token $x'_{m}$ within the ground truth mask $M_{gt}$. These probabilities are computed using a 2D Gaussian distribution, where the distribution parameters are generated based on the reference token $x_{m}$. In other words, we aim to ensure that each token $x_{m}$ within the ground truth mask $M_{gt}$ can generate a 2D Gaussian distribution that effectively covers the intended IOI.

In practice, obtaining an accurate Gaussian distribution that closely mimic the ground-truth mask poses a granular alignment challenge. For example, when the gaze point lands on the headlight of a car, it is often ambiguous whether the target object is the headlight itself or the entire vehicle. Similarly, if the gaze fixed on a chess piece, it is unclear whether the segmentation should cover the individual piece or the whole chessboard. To mitigate this ambiguity, we introduce an additional classification loss \(\mathcal{L}_{\mathrm{cls}}\), which leverages category labels to enforce granular alignment. Beyond alignment, \(\mathcal{L}_{\mathrm{cls}}\) also provide clustering effect, markedly enhancing the inter‑class diversity of the Gaussian maps. As an auxiliary supervision signal, it encourages the saliency encoder $S(.)$ to learn class‑specific features, leading to a better Gaussian maps. Therefore, as shown in Figure~\ref{fig:fosam_framework}, a shallow neural network $U(.)$ is attached to the output of $S(.)$ for classification task. Let $M_{cls}(x_{i})$ denote the class label associated with the IOI where $x_{m}$ is mapping to, $\mathcal{L}_{\mathrm{cls}}$ can be expressed as:
\begin{equation}
\label{eqn:loss2}
\small
\mathcal{L}_{\mathrm{cls}} = -\sum_{x_{m}\in X_{M_{gt}}} CE[U(S(x_{m}),M_{cls}(x_{m})]
\end{equation}
Here, $CE(\cdot)$ denotes the cross-entropy loss. Finally, to ensure that the output $M_{pred}$ of the image decoder $D(\cdot)$ aligns with the ground-truth mask $M_{gt}$, we apply the Dice loss.
\begin{equation}
\label{eqn:loss3}
\small
\mathcal{L}_{\mathrm{dice}} = DICE[M_{pred},M_{gt}]
\end{equation}
The overall loss function can be calculated as $\mathcal{L}_{\mathrm{total}} = \lambda_{1} \mathcal{L}_{\mathrm{nll}}+ \lambda_{2} \mathcal{L}_{\mathrm{cls}} + \lambda_{3} \mathcal{L}_{\mathrm{dice}}$, where $\lambda_{1}$,$\lambda_{2}$ and $\lambda_{3}$ represent the relative importance of the loss functions. During training, only the saliency encoder $S(\cdot)$, classification head $U(\cdot)$, Gaussian head $G(\cdot)$, and image decoder $D(\cdot)$ are updated, while the image encoder $E(\cdot)$ remains frozen, significantly reducing the training cost of FoSAM.

\begin{comment}
As indicated in Figure~\ref{fig:fosam_fig}, the inference‑only path is shown in blue, while the training‑only path is shown in red. Unlike inference, where only the gaze token \(\bm{\tilde{z}}_{g}\) is forwarded, the training phase feeds all tokens covered by the ground‑truth mask \(\bm{M}\in\{0,1\}^{H^{I}\times W^{I}}\) to the subsequent network. The ground‑truth mask is first down‑sampled to \(\tilde{H}^{t}\times\tilde{W}^{t}\), yielding \(\bm{\tilde{M}}\in\mathbb{R}^{\tilde{H}^{t}\times\tilde{W}^{t}}\).
Then the resulting set of masked tokens can be defined as

\begin{equation}
\bigl\{\bm{\tilde{z}}_{i}\bigr\}^{m}
   = \Bigl\{\bm{\tilde{z}}_{i}\;\Bigl|\;
      \bm{\tilde{M}}(h_{i},w_{i})>0\Bigr\}.
\end{equation}


The motivation is intuitive: all masked tokens make up parts of the same object, and during inference a gaze point landing on any token in \(\{\bm{\tilde{z}}_{i}\}^{m}\) should predict the same Gaussian. Consequently, using every masked token as a training sample effectively enhances the model’s generalization capability (\textcolor{red}{at this point, we have not define the training loss, it is very strange to talk about any intuition}).

During training, we need to optimize three loss terms: the negative log‑likelihood
\(\mathcal{L}_{\mathrm{nll}}\), the classification loss \(\mathcal{L}_{\mathrm{cls}}\), and the Dice loss \(\mathcal{L}_{\mathrm{dice}}\).

\noindent\textbf{Negative log-likelihood Loss.}
The purpose of \(\mathcal{L}_{\mathrm{nll}}\) is to teach the Gaussian predictor to predict a 2D Gaussian distribution from the masked token \(\{\bm{\tilde{z}}_{i}\}^m\). It can be expressed as:
\begin{equation}
\begin{aligned}
\mathcal{L}_{\mathrm{nll}}&=
  -\frac{1}{N^2}\sum_{j=1}^{N}\sum_{i=1}^{N}
     \log G_{\theta_j}(h_i, w_i)\bm{\tilde{M}}(h_i,w_i),\\
\theta_j &= \bigl[\bm{\mu},\bm{\sigma},\rho \bigr]=g^{\text{Gau}}(\bm{\tilde{z}}_{j}\in \bigl\{\bm{\tilde{z}}_{i}\bigr\}^{m})\\
N &= \bigl|\bigl\{\bm{\tilde{z}}_{i}\bigr\}^{m}\bigr|,\\
\end{aligned}
\end{equation}

where \((h_i, w_i)\) denotes the coords of the token \(\bm{\tilde{z}}_{i} \in \{\bm{\tilde{z}}_{i}\}^{m}\), \(\theta_j\) denotes the Gaussian parameters and \(g^{\text{Gau}}\) is the Gaussian predictor. This loss treats every pixel in the ground‑truth mask \(\bm{\tilde{M}}\) as samples and measures the distributional similarity between the predicted Gaussian and \(\bm{\tilde{M}}\).

\noindent\textbf{Classification Loss.} In practice, obtaining an accurate Gaussian window poses a granular alignment challenge. For example, when the gaze point lands on the headlight of a car, it is often ambiguous whether the target object is the headlight itself or the entire vehicle. Similarly, if the gaze fixed on a chess piece, it is unclear whether the segmentation should cover the individual piece or the whole chessboard.  

To mitigate this ambiguity, we introduce an additional classification loss \(\mathcal{L}_{\mathrm{cls}}\), which leverages category labels to enforce granular alignment. Beyond alignment, \(\mathcal{L}_{\mathrm{cls}}\) also provide clustering effect, markedly enhancing the inter‑class diversity of the Gaussian windows. As an auxiliary supervision signal, it encourage the saliency encoder to embed category‑specific features into the token representations \(\{\bm{\tilde{z}}_i\}\), thereby improving the estimation of the Gaussian window’s shape and scale. \(\mathcal{L}_{\mathrm{cls}}\) can be expressed as:

\begin{equation}
\begin{aligned}
\mathcal{L}_{\mathrm{cls}} &= \frac{1}{N}\sum_{i=1}^{N}\operatorname{CE}\bigl(g^{\text{cls}}(\bm{\tilde{z}}_i\in\{\bm{\tilde{z}}_i\}^m),~\text{Label}\bigr),\\
N &= \bigl|\bigl\{\bm{\tilde{z}}_{i}\bigr\}^{m}\bigr|,\\
\end{aligned}
\end{equation}

where \(\operatorname{CE}\) is Cross Entropy and \(g^{\text{cls}}\) is classification head.
\end{comment}

\subsection{FoSAM Streaming Algorithm}
\label{sec:method:frame_skip}

\begin{wrapfigure}{r}{0.25\textwidth}
  \begin{center}
    \includegraphics[width=0.25\textwidth]{figures/skip_algo.pdf}
  \end{center}
  \caption{FSA flow.}
  \label{fig:skip_algo}
  \vspace{-10pt}
\end{wrapfigure}

Building on the FoSAM architecture introduced in Section~\ref{sec:fosam_frameowrk}, this section explains how FoSAM can be extended for efficient instance segmentation across consecutive video frames. As motivated by the study in Section~\ref{sec:bg:eye-behavior}, FoSAM is triggered during real-time operation only when the input image $I$, captured by the front-facing camera, shows a significant change or when a shift in the user's gaze direction is detected. The logic flow of the FoSAM Streaming Algorithm (FSA) is shown in Figure~\ref{fig:skip_algo}.

 
Let $I_{ini}$ and $I_{t}$ represent the input images at the initial frame and at time step $t$, respectively. A simple criterion is used to determine whether the current view remains within the previous segment by computing the difference between $I_{t}$ and $I_{\text{ini}}$ (Condition 1). If this difference surpasses a predefined threshold $\alpha$, it signals a substantial change in the front-facing scene, prompting a full re-execution of FoSAM, and $I_{ini}$ will be updated with $I_{t}$. If not big change is observed, the current gaze location $g_t$ is examined to determine whether it remains in IOI region defined by the segmentation mask $M_{pred}$ of $I_{ini}$ (Condition 2). If it does, $M_{pred}$ can be reused; otherwise FoSAM must be executed with $g_t$ and $M_{pred}$ will be updated accordingly. A detailed description of the FSA is provided in the supplementary materials.

\section{Evaluation}
\begin{wraptable}{r}{0.6\textwidth}
    \centering\resizebox{0.6\columnwidth}{!}{
       \begin{tabular}{c|c|c|c|c|c}
        \toprule
        Method & K   & ADE20K & LVIS & Cityscape & GFLOP \\ \midrule
        
        SAM-B~\cite{ravi2024sam} & 4096 & 0.511 & 0.537   & 0.347 & 744\\  
        %ESAM-S~\cite{xiong2024efficientsam} & 1600 & $640 \times 640$ & \textcolor{red}{0.330}& 73.6 & \textcolor{red}{0.248} & 73.6 & \textcolor{red}{0.321} & 73.6\\ 
        %ESAM-T~\cite{xiong2024efficientsam} & 1600 & $640 \times 640$ & \textcolor{red}{0.250}& \textcolor{red}{21.8} & \textcolor{red}{0.240}& \textcolor{red}{21.8} & \textcolor{red}{0.210}& \textcolor{red}{21.8}\\ 
        ESAM-S~\cite{xiong2024efficientsam} & 4096  & 0.411  & 0.359  & 0.304 & 188.2\\ 
        ESAM-T~\cite{xiong2024efficientsam} & 4096   & 0.350   & 0.236 & 0.248 & 56.1 \\ 
        SlimSAM-S~\cite{Chen2024NEURIPS} & 4096   & 0.477  & 0.474 & 0.293 & 15.1 \\
        SlimSAM-T~\cite{Chen2024NEURIPS} & 4096 & 0.461 & 0.458 & 0.263 & 11.0\\
        FSNet-DL~\cite{zeng2025foveated} & NA  & 0.34  & 0.38  & 0.24 & 18.8 \\ 
        FSNet-SF~\cite{zeng2025foveated} & NA  & 0.33  & 0.36  & 0.21 & 12.6\\ \midrule

        AD-S & 200  & 0.273  & 0.350  & 0.221 & 14.6\\
        LC-S & 200  & 0.368  & 0.477  & 0.385 & 14.6\\
        \rowcolor{blue!10}
        \textbf{FoSAM-S} & 200   & \textbf{0.495} & \textbf{0.487}  & \textbf{0.404} & \textbf{14.6}\\ \midrule
          
        AD-S & 100 & 0.252  & 0.343  & 0.203 & 10.4\\ 
        LC-S & 100  & 0.294  & 0.444  & 0.339 & 10.4\\ 
        \rowcolor{blue!10}
        \textbf{FoSAM-S} & 100   & \textbf{0.465}  & \textbf{0.461}  & \textbf{0.362} & \textbf{10.4}\\ \midrule          

        AD-T & 200  & 0.256  & 0.322  & 0.198 & 7.4\\
        LC-T & 200  & 0.327  & 0.453  & 0.371 & 7.4\\
        \rowcolor{blue!10}
        \textbf{FoSAM-T} & 200  & \textbf{0.471}  & \textbf{0.469}  & \textbf{0.377} & \textbf{7.4}\\ \midrule
          
        AD-T & 100 & 0.239  & 0.318  & 0.176 & 5.3\\           
        LC-T & 100  & 0.264  & 0.421  & 0.302 & 5.3\\ 
        \rowcolor{blue!10}
        \textbf{FoSAM-T} & 100 & \textbf{0.457}  & \textbf{0.435}  & \textbf{0.328} & \textbf{5.3}\\ 
        \bottomrule
        \end{tabular}
}
\caption{IoU and FLOP comparison of FoSAM and baseline algorithms. K denotes the token budget.}
\label{tab:eval:main}
\vspace{-5pt}
\end{wraptable}

We evaluate our approach using three publicly available datasets: ADE20K~\cite{zhou2019ade20k}, LVIS~\cite{gupta2019lvisdatasetlargevocabulary}, and Cityscapes~\cite{Cordts2016Cityscapes}. As these datasets are originally designed for generic segmentation tasks to cover all the objects in the image, we introduce a gaze-aware masking preprocessing step to enable segmentation over the IOI only. Specifically, for each training and testing sample, a gaze location is randomly selected within the image, and the corresponding IOI region is defined based on this location. The size of the input images is $640\times 640$ for all datasets. %The training and test sets contain 60,000 and 12,000 samples, respectively, with all input images resized to $640 \times 640$.

We use a lightweight two-layer ViT as the saliency encoder $S(.)$. The input images have a resolution of $H \times W = 640 \times 640$, and are average downsampled to $H'_e \times W'_e = 160 \times 160$ prior to processing. The FoSAM framework is integrated with ESAM Small~\cite{xiong2024efficientsam} and ESAM Tiny, referred to as FoSAM-S and FoSAM-T, respectively. Two baseline methods are developed for comparison. The first method,~\textit{Average Downsampling} (AD), enhances segmentation efficiency by directly applying average pooling to the original input $X$, reducing its spatial resolution and thus lowering computational cost. The second method,~\textit{Local Cropping (LC)}, improves efficiency by extracting a fixed-size patch centered around the gaze point $g$, thereby narrowing the focus to a smaller input region. These baselines are evaluated on both ESAM Small and ESAM Tiny, resulting in AD-S, AD-T, LC-S, and LC-T configurations. Furthermore, we evaluate performance against SAM Base (SAM-B), ESAM Small (ESAM-S), ESAM Tiny (ESAM-T), and SlimSAM~\cite{Chen2024NEURIPS}, which improves segmentation efficiency through a combination of knowledge distillation, embedding pruning, and bottleneck pruning. We also compare FoSAM with FSNet, which integrates DeepLab~\cite{he2016deep} and Segformer~\cite{xie2021segformer} backbones to enable efficient segmentation via learnable input downsampling strategies, referred to as FSNet-DL and FSNet-SF, respectively. 
During training, the weights for the loss terms $\mathcal{L}_{\mathrm{nll}}$, $\mathcal{L}_{\mathrm{cls}}$, and $\mathcal{L}_{\mathrm{dice}}$ are all set to 1. We evaluate segmentation performance using the Intersection over Union (IoU) metric, which measures the overlap between the predicted segmentation and the ground truth IOI region. To control the average number of retained tokens $K$, we adjust the saliency score threshold $\tau$ and the minimum token count $\epsilon$ as defined in Equation~\ref{eqn:thres}. More evaluation results are shown in the supplementary materials.


\begin{comment}
\begin{table}
    \centering\resizebox{\columnwidth}{!}{
       \begin{tabular}{c|c|c|cc|cc|cc}
        \toprule
        \multirow{2}{*}{Method} & \multirow{2}{*}{Token Budget (K)}  & \multirow{2}{*}{Image Size} & \multicolumn{2}{c|}{ADE20K} & \multicolumn{2}{c|}{LVIS} & \multicolumn{2}{c}{Cityscape} \\
        & & & IoU & GFLOP & IoU & GFLOP & IoU & GFLOP  \\ \midrule
        
        SAM-B~\cite{ravi2024sam} & 4096 & $1024 \times 1024$ & 0.511 & 744& 0.537 & 744  & 0.347 & 744\\  
        %ESAM-S~\cite{xiong2024efficientsam} & 1600 & $640 \times 640$ & \textcolor{red}{0.330}& 73.6 & \textcolor{red}{0.248} & 73.6 & \textcolor{red}{0.321} & 73.6\\ 
        %ESAM-T~\cite{xiong2024efficientsam} & 1600 & $640 \times 640$ & \textcolor{red}{0.250}& \textcolor{red}{21.8} & \textcolor{red}{0.240}& \textcolor{red}{21.8} & \textcolor{red}{0.210}& \textcolor{red}{21.8}\\ 
        ESAM-S~\cite{xiong2024efficientsam} & 4096 & $1024 \times 1024$ & 0.411 & 188.2 & 0.359 & 188.2 & 0.304 & 188.2\\ 
        ESAM-T~\cite{xiong2024efficientsam} & 4096 & $1024 \times 1024$  & 0.350 & 56.1  & 0.236 & 56.1& 0.248 & 56.1 \\ 
        SlimSAM-S~\cite{Chen2024NEURIPS} & 4096 & $1024 \times 1024$  & 0.443 & 15.1 & 0.474 & 15.1& 0.293 & 15.1 \\
        %SlimSAM-T~\cite{Chen2024NEURIPS} & 4096 & $1024 \times 1024$  & 0.424 & 11.0 & 0.458 & 11.0 & 0.262 & 11.0\\
        FSNet-DL~\cite{zeng2025foveated} & NA & $80 \times 80$ & 0.34 & 18.8 & 0.38 & 18.8 & 0.24 & 18.8 \\ 
        FSNet-SF~\cite{zeng2025foveated} & NA & $80 \times 80$ & 0.33 & 12.6 & 0.36 & 12.6 & 0.21 & 12.6\\ \midrule

        AD-S & 200 & $640 \times 640$ & 0.273 & 14.6 & 0.350 & 14.6 & 0.221 & 14.6\\
        LC-S & 200 & $640 \times 640$ & 0.368 & 14.6 & 0.477 & 14.6 & 0.385 & 14.6\\
        \rowcolor{blue!10}
        \textbf{FoSAM-S} & 200 & $640 \times 640$  & \textbf{0.495} & \textbf{14.6}& \textbf{0.487} & \textbf{14.6} & \textbf{0.404} & \textbf{14.6}\\ \midrule
          
        AD-S & 100 & $640 \times 640$ & 0.252 & 10.4 & 0.343 & 10.4 & 0.203 & 10.4\\ 
        LC-S & 100 & $640 \times 640$ & 0.294 & 10.4 & 0.444 & 10.4 & 0.339 & 10.4\\ 
        \rowcolor{blue!10}
        \textbf{FoSAM-S} & 100 & $640 \times 640$  & \textbf{0.465} & \textbf{10.4} & \textbf{0.461} & \textbf{10.4} & \textbf{0.362} & \textbf{10.4}\\ \midrule          

        AD-T & 200 & $640 \times 640$ & 0.256 & 7.4 & 0.322 & 7.4 & 0.198 & 7.4\\
        LC-T & 200 & $640 \times 640$ & 0.327 & 7.4 & 0.453 & 7.4 & 0.371 & 7.4\\
        \rowcolor{blue!10}
        \textbf{FoSAM-T} & 200 & $640 \times 640$ & \textbf{0.471} & \textbf{7.4} & \textbf{0.469} & \textbf{7.4} & \textbf{0.377} & \textbf{7.4}\\ \midrule
          
        AD-T & 100 & $640 \times 640$ & 0.239 & 5.3 & 0.318 & 5.3 & 0.176 & 5.3\\           
        LC-T & 100 & $640 \times 640$ & 0.264 & 5.3 & 0.421 & 5.3 & 0.302 & 5.3\\ 
        \rowcolor{blue!10}
        \textbf{FoSAM-T} & 100 & $640 \times 640$ & \textbf{0.457} & \textbf{5.3} & \textbf{0.435} & \textbf{5.3} & \textbf{0.328} & \textbf{5.3}\\ 
        \bottomrule
        \end{tabular}
}
\caption{IoU and FLOP comparison of FoSAM and baseline algorithms.}
\label{tab:eval:main}
\end{table}    
\end{comment}


\subsection{Evaluation Results of FoSAM}

\begin{figure}[t]
  \centering
  \begin{minipage}{0.25\textwidth}
    \centering
    \vspace{1pt}
    \resizebox{\textwidth}{!}{
      \begin{tabular}{c|c c}
        \toprule
        Latency  & \multicolumn{2}{c}{Token} \\
        (ms) & 200 & 100 \\ \midrule
        FoSAM-S & 21.6  & 18.3  \\
        FoSAM-T & 17.6  & 14.5  \\
        \bottomrule
      \end{tabular}
    }
    \captionof{table}{Processing latencies of FoSAM on an AR platform.}
    \label{tab:fosam_latency}
  \end{minipage}%
  \hspace{0.5em} % 控制两者间距，可调大调小
  \begin{minipage}{0.5\textwidth}
    \centering
    \includegraphics[width=0.9\textwidth]{figures/skip_eval.pdf}
    \caption{Average IoU and frame skip rate of FSA on FoSAM-S and FoSAM-T under varying $\alpha$. }
    \label{fig:skip_eval}
  \end{minipage}
  \hspace{0.5em} % 控制两者间距，可调大调小
  \begin{minipage}{0.2\textwidth}
    \centering
    \includegraphics[width=0.8\textwidth]{figures/userstudy.jpg}
    \caption{Participants are doing user study. }
    \label{fig:userstudy}
  \end{minipage}
\end{figure}


%  \begin{wrapfigure}{r}{0.6\textwidth}
%   \begin{center}
%     \includegraphics[width=0.6\textwidth]{figures/skip_eval.pdf}
%   \end{center}
%   \caption{Average IoU of FSA under varying $\alpha$ and $\beta$ settings using (a) the FoSAM-S segmentation model and (b) the FoSAM-T segmentation model.}
%   \label{fig:skip_eval}
% \end{wrapfigure}
%\textcolor{red}{more explanation}
Table~\ref{tab:eval:main} presents the evaluation results. The compared methods, FoSAM, AD, and LC, all operate on input images at their original resolution. FoSAM adopts an adaptive token sampling strategy with average token budgets of $K = 200$ and $K = 100$ across the entire test dataset. To ensure a fair comparison, the input sizes of AD and LC are adjusted so that their computational cost aligns with that of FoSAM. For ESAM, SlimSAM, and SAM, regardless of the original input size, the built-in resizing layer by default resizes all inputs to $1024 \times 1024$, resulting in a fixed number of 4096 tokens. For SlimSAM, we tune its pruning ratio to achieve a similar computational cost to FoSAM. In the case of FSNet-DL and FSNet-SF, input images are downsampled to a resolution of $80 \times 80$, which is the default setting of it. As shown in the results, FoSAM consistently outperforms the AD and LC baselines across all datasets. In contrast to different versions of ESAM, SlimSAM and FSNet, FoSAM achieves both higher accuracy and reduced computation. In the ADE20K and Cityscape datasets, FoSAM-T with 100 token budget achieves $1.4\%$ and  $3.5\%$ higher IoU than SlimSAM with only one third of the computation cost. This highlights the effectiveness of IOI guided adaptive token sampling in FoSAM. Moreover, FoSAM demonstrates strong generalizability across datasets, with minor performance variations attributed to dataset-specific characteristics. Compared to SAM, although FoSAM does not achieve equally high IoU, it provides comparable accuracy while reducing computational cost by over $50\times$. As shown in the user study in Section~\ref{sec:user-study}, FoSAM also offers a significantly better visual experience, whereas SAM introduces noticeable processing delays that negatively impact user experience.


%\paragraph{Execution Latency of FoSAM on AR Device} 
To assess FoSAM's processing latency on AR device, we profile the performance of FoSAM-S and FoSAM-T with average token budgets of K=200 and 100 on the Qualcomm XR2 Gen2 platform~\cite{qualcomm_xr2_gen2}, utilizing the Qualcomm AI Hub toolkit~\cite{qualcomm_aihub}.
It is important to note that the Qualcomm XR2 Gen2 is integrated into the Meta Quest Pro for AR/VR task processing~\cite{quest_pro, engel2023project}. The latency results on a single LVIS sample are reported in Table~\ref{tab:fosam_latency}. Compared to the latency of SAM and ESAM shown in Table~\ref{tab:reduce_size}, FoSAM achieves an average of $24\times$ reduction in inference time and meets the 30 ms latency threshold required for a seamless visual experience, as indicated by prior research~\cite{kaaresoja2014towards, visual_requirement,albert2017latency}.


\subsection{Evaluation Results of FoSAM Streaming Algorithm}
The FSA mechanism introduced in Section~\ref{sec:method:frame_skip} reduces segmentation overhead by reusing results from previous frames. However, as illustrated in Figure~\ref{fig:skip_algo}, varying the parameters can influence the average IoU across frames. To evaluate this effect, we analyze the performance of FoSAM-S and FoSAM-T with FSA evaluated on the Cityscapes dataset under different settings of $\alpha$, which represent the threshold for image difference shown in Figure~\ref{fig:skip_algo}. The token budget is set to 200 for both FoSAM-S and FoSAM-T. We use 10 video sequences which contain 5000 images in total. However, since the Cityscapes dataset lacks gaze location data, we incorporate gaze traces described in~\cite{lv2024ariaeverydayactivitiesdataset} into each frame of the Cityscapes dataset. The results are described in Figure~\ref{fig:skip_eval}, as $\alpha$ increases from 0.01 to 0.1, more frames are skipped by reusing segmentation results of previous frames.
For instance, skipping $40\%$ of the frames leads to only a 0.03 reduction in average IoU across frames for both FoSAM-S and FoSAM-T, showing that FSA can greatly reduce segmentation computation with negligible impact on the accuracy. 

\subsection{Ablation Study}
\paragraph{Impact of the Token Budget}

  \begin{wraptable}{r}{0.35\textwidth}
    \centering\resizebox{0.35\columnwidth}{!}{
       \begin{tabular}{c|c|cc}
        \toprule
        Dataset & Token Budget & FoSAM-S & LC-S   \\ \midrule
        \multirow{3}{*}{LVIS} & 400 & 0.499 & 0.490 \\
        & 200 & 0.487 & 0.477\\
        & 100 & 0.475 & 0.443 \\
        & 25 & 0.371 & 0.318 \\ \midrule
        \multirow{3}{*}{Cityscape} & 400 & 0.417 & 0.394 \\
        & 200 & 0.404 & 0.385 \\
        & 100 & 0.362 & 0.339 \\
        & 25 & 0.269 & 0.243\\ \midrule
        \multirow{3}{*}{ADE} & 400 & 0.517 & 0.442 \\
        & 200 & 0.499 & 0.368 \\
        & 100 & 0.472 & 0.293\\
        & 25 & 0.311 & 0.227 \\
        
        \bottomrule
        \end{tabular}
}
\caption{FoSAM and LC performances with varying token budgets. }
\label{tab:ablation:window_size}
\end{wraptable}
\begin{comment}
  \begin{wraptable}{r}{0.5\textwidth}
    \centering\resizebox{0.5\columnwidth}{!}{
       \begin{tabular}{c|c|cc|cc}
        \toprule
        Dataset & Token Budget & FoSAM-S & LC-S & FoSAM-T & LC-T  \\ \midrule
        \multirow{3}{*}{LVIS} & 400 & 0.500 & 0.490 & - & -\\
        & 200 & 0.487 & 0.477 & - & -\\
        & 100 & 0.475 & 0.443 & - & -\\
        & 25 & 0.371 & 0.318 & - & -\\ \midrule
        \multirow{3}{*}{Cityscape} & 400 & 0.417 & 0.394 & 0.394 & 0.372\\
        & 200 & 0.404 & 0.385 & 0.377 & 0.371\\
        & 100 & 0.362 & 0.339 & 0.328 & 0.302\\
        & 25 & 0.269 & 0.243 & 0.225 & 0.199\\ \midrule
        \multirow{3}{*}{ADE} & 400 & 0.517 & 0.442 & 0.495 & 0.418\\
        & 200 & 0.499 & 0.368 & 0.471 & 0.327\\
        & 100 & 0.472 & 0.293 & 0.457 & 0.264\\
        & 25 & 0.311 & 0.227 & 0.294 & 0.205\\
        
        \bottomrule
        \end{tabular}
}
\caption{Comparison between FoSAM and Local Cropping across different token budgets. }
\label{tab:ablation:window_size}
\end{wraptable}
\end{comment}

As described in Section~\ref{sec:dynamic-window}, FoSAM adopts an adaptive token sampling mechanism that adjusts the number of retained tokens according to the size of the IOI. In contrast, the LC algorithm applies a fixed token count regardless of IOI size, which can lead to too few tokens for large IOIs, reducing accuracy, or too many tokens for small IOIs, causing unnecessary computation. Table~\ref{tab:ablation:window_size} compares the performance of FoSAM and LC under varying token budgets $K$ on different datasets. For FoSAM, the number of retained tokens varies across frames, with an average of $K$. On the LVIS dataset, FoSAM-S achieves an average of $3.6\%$ higher IoU than LC-S when $K = 400$, and outperforms LC-S by $5.5\%$ in IoU when $K = 25$. These results highlight FoSAM’s ability to adapt the token budget to the size of the IOI, allowing it to maintain accuracy even under constrained token budgets.

\begin{comment}  
% move this to the supplementary materials
\subsubsection{Impact of Adaptive Token Sampling Hyperparameters}

\begin{wraptable}{r}{0.3\textwidth}
    \centering\resizebox{0.3\columnwidth}{!}{
       \begin{tabular}{c|c|c|c}
        \toprule
        $\tau$ & $c$ & FoSAM-S & Token Budget  \\ \midrule
        0.2 & 300 & 0.404 & 400 \\
        0.12 & 200 & 0.417 & 400\\
        0.1 & 100 & 0.416 & 400\\ \midrule
        0.4 & 150 & 0.39 & 200\\
        0.3 & 100 & 0.404 & 200\\
        0.25 & 50 & 0.402 & 200\\ \midrule
        0.7 & 75 & 0.355 & 100\\
        0.6 & 50 & 0.362 & 100\\
        0.55 & 25 & 0.362 & 100\\
        \bottomrule
        \end{tabular}
}
\caption{Evaluation results of adaptive token sampling settings.}
\label{tab:ablation:hyper}
\end{wraptable}

The saliency score threshold \(\tau\) and the minimum token count \(c\) are used to determine how many tokens are retained, as detailed in Section~\ref{sec:dynamic-window}. Different combinations of \(\tau\) and \(c\) can result in the same average number of tokens being retained across the samples in a dataset. In Table~\ref{tab:ablation:hyper}, we evaluate the effects of varying \(\tau\) and \(c\) under different token budgets using FoSAM-S over the Cityscape dataset. A larger \(\tau\) reduces the number of tokens selected based on saliency scores; to maintain the target token budget, \(c\) must be increased accordingly. However, a larger \(c\) restricts FoSAM’s ability to allocate fewer tokens for small IOIs. For example, under a token budget of 200, the configuration \(\tau = 0.3, c = 100\) achieves better performance than other settings.


\subsubsection{Impact of Loss Weights}

During the training of FoSAM, the overall loss is computed as a weighted sum of the negative log-likelihood loss, classification loss, and Dice loss, with weights \(\lambda_1\), \(\lambda_2\), and \(\lambda_3\), respectively, as defined in Equation~\ref{eqn:loss1}, Equation~\ref{eqn:loss2}, and Equation~\ref{eqn:loss3}. Dice loss plays a critical role in the segmentation task, while the negative log-likelihood loss and classification loss help the saliency encoder in FoSAM better adapt the saliency score distribution based on the IOI. In Table~\ref{tab:ablation:loss}, we evaluate the contributions of the negative log-likelihood loss and the classification loss to the overall performance of FoSAM by varying their weights \(\lambda_1\) and \(\lambda_2\), using the FoSAM-S model on the Cityscapes dataset. 

\begin{wraptable}{r}{0.3\textwidth}
    \centering\resizebox{0.3\columnwidth}{!}{
       \begin{tabular}{c|c|c|c}
        \toprule
        $\lambda_1$ & $\lambda_2$ & Token Budget & IoU  \\ \midrule
        1 & 1 & 200 & 0.404 \\
        2 & 1 & 200 & 0.398 \\
        1 & 0 & 200 & 0.381 \\

        \bottomrule
        \end{tabular}
}
\caption{Impact of the weights of the nll loss ($\lambda_1$) and classification loss ($\lambda_2$).}
\label{tab:ablation:loss}
\end{wraptable}

The setting \(\lambda_1 = 1, \lambda_2 = 1\) is used as the baseline for our main evaluation results. Increasing \(\lambda_1\) reduces the relative contribution of the classification loss, which leads to a decrease in segmentation accuracy. When \(\lambda_2 = 0\), the classification loss is completely removed, resulting in a further drop in accuracy. This highlights the importance of the classification loss in achieving good segmentation performance.
\end{comment}
 

\paragraph{Impact of Saliency Encoder Size}

\begin{wraptable}{r}{0.4\textwidth}
    \centering\resizebox{0.4\columnwidth}{!}{
       \begin{tabular}{c|c|c|c}
        \toprule
        Layer Number & Image Size & IoU & SE GFLOPs  \\ \midrule
        3 & $640 \times 640$ & 0.499 & 10.4 \\
        3 & $320 \times 320$ & 0.499 & 2.8 \\ \midrule
        2 & $320 \times 320$ & 0.497 & 1.88 \\
        2 & $160 \times 160$ & 0.495 & 0.48 \\
        2 & $80 \times 80$ & 0.492 & 0.12 \\ \midrule
        1 & $320 \times 320$ & 0.492 & 0.94 \\

        \bottomrule
        \end{tabular}
}
\caption{Impact of the saliency encoder.}
\label{tab:ablation:layer}
\vspace{-10pt}
\end{wraptable}
In this section, we examine how the computational cost of the saliency encoder $S(\cdot)$ influences FoSAM's accuracy under an average token budget of $K = 200$ on the ADE20K dataset. Specifically, we evaluate the impact on FoSAM performance by varying both the parameter size of $S(\cdot)$ and the input resolution of $X'$ fed into $S(\cdot)$.
As shown in Table~\ref{tab:ablation:layer}, increasing the number of encoder layers from 2 to 3 and raising the input resolution from $160 \times 160$ to $640 \times 640$ leads to only a slight improvement in IoU, while the computational cost rises substantially from 0.48 GFLOPs to 10.4 GFLOPs for the saliency encoder processing. Additionally, reducing the number of encoder layers results in a more significant accuracy drop than downsampling the input image. The saliency encoder design in FoSAM strikes an effective balance between computational efficiency and segmentation performance.

\subsection{User Study}
\label{sec:user-study}


% \begin{wraptable}{r}{0.25\textwidth}
%     \centering\resizebox{0.25\columnwidth}{!}{
%        \begin{tabular}{c|c}
%         \toprule
%     FoSAM-S & FoSAM-T \\ \midrule
%     21.6 ms & 17.6 ms\\
%     \bottomrule
%     \end{tabular}
% }

% \caption{Processing latencies of FoSAM on an edge GPU.}
% \label{tab:reduce_size}
% % \vspace*{-1\baselineskip}
% \end{wraptable}

\begin{wrapfigure}{r}{0.22\textwidth}
\vspace{-15pt}
\centering \includegraphics[width=0.22\textwidth]{figures/userstudy_result.pdf}
\caption{Preference rates for FoSAM among the seven participants.}
\label{fig:userstudy-results}
\end{wrapfigure}
To assess the user experience benefits of FoSAM over the standard SAM method on VR hardware, we conducted a two-interval forced-choice (2IFC) user study simulating both methods’ effects on the Meta Quest Pro headset~\cite{quest_pro}. Segmentation masks were precomputed and gaze-contingent visualizations were rendered with artificially introduced latencies reflecting each method’s runtime profiled on the Qualcomm XR2 Gen2's NPU~\cite{qualcomm_hexagon_npu}. Each of the seven participants compared results across four test images, selecting the preferred output in 32 trials, as shown in Figure~\ref{fig:userstudy}. As shown in Figure~\ref{fig:userstudy-results}, FoSAM was preferred in 96.9\%$\pm$4.8\% of trials overall, consistently outperforming SAM across all images (98.2\%$\pm$4.7\% for Image 1, 96.4\%$\pm$6.1\% for Image 2, 96.4\%$\pm$9.4\% for Image 3, and 96.4\%$\pm$9.4\% for Image 4). Although FoSAM incurs a slight drop in segmentation accuracy compared to SAM, its substantially lower latency and better temporal alignment with gaze location result in a significantly improved user experience. These results validate the effectiveness of our latency-accuracy trade-off design for AR/VR applications. See supplementary materials for implementation details.

\section{Conclusion}
\label{sec:conclusion}
We presented FoSAM, a gaze-guided segmentation framework optimized for AR/VR. By focusing computation on user-relevant regions, FoSAM achieves over $50\times$ speedup with minimal accuracy loss. Its extension, FSA further improves efficiency through result reuse across stable frames. Experiments and user studies validate its effectiveness for real-time deployment on resource-limited devices.






{
\small
\newpage
\bibliographystyle{plain}
\bibliography{ref}

}


%%%%%%%%%%%%%%%%%%%%%%%%%%%%%%%%%%%%%%%%%%%%%%%%%%%%%%%%%%%%
\begin{comment}
\appendix
\newpage
\section{Technical Appendices and Supplementary Material}
Technical appendices with additional results, figures, graphs and proofs may be submitted with the paper submission before the full submission deadline (see above), or as a separate PDF in the ZIP file below before the supplementary material deadline. There is no page limit for the technical appendices.


    

\subsection{FoSAM Streaming Algorithm}
\textcolor{blue}{TODO}

\begin{algorithm}[t]
\caption{FoSAM-Stream Algorithm}\label{alg:reuse}
\small
\KwIn{$I_{f}^{l}$ is the initial frame of current VS. $g^{l}$ and $Y^{last}$ are buffered gaze location and segmentation label map. \(g^t\), \(I^{t}_{f}\), \(I^{dt}_{f}\), and \(Y_{cm}^{t}\) represent the current gaze direction, the input frame, the corresponding uniformly subsampled version of it, and the segmentation label map at time \(t\), $SAC$ is the saccade detection signal. $\alpha$ and $\beta$ are the thresholds of result reuse. G(.) is the gaze tracking function.}
\SetKwBlock{Begin}{Initiation}{}
\Begin{
    $I_{f}^{l} = \varnothing$, $g^{l} = \varnothing$, $Y^{last} = \varnothing$\\
    \For{$1\leq t \leq T$}{
    \If{$\sum_{ij}|I^{dt}_{f,ij}-I^{l}_{f,ij}| > \alpha$}{
            Run XXXNet with $I^{t}_{f}$ and $I^{t}_{e}$, get $Y_{cm}^{t}$\;
            $I_{f}^{l} \gets I^{dt}_{f}$, $g^{l} \gets G(I^{t}_{e})$, $Y^{last} \gets Y_{cm}^{t}$\;
            \KwRet $Y_{cm}^{t}$
        }
    \Else{
        \If{$SAC==1$}{
        Saccade detect, halt rest operations.
    }
        \Else{
            $g^t = G(I^{t}_{e})$ \\
            \If{$|g^t-g^{l}|^{2} > \beta$}{
                Run XXXNet with $I^{t}_{f}$ and $I_{e}$, get $Y_{cm}^{t}$\;
                $g^{l} \gets g^t$, $Y^{last} \gets Y_{cm}^{t}$\;
                \KwRet $Y_{cm}^{t}$
            }
            \Else{
                \KwRet $Y^{last}$
            }
        }
    }
    }
}
\end{algorithm}

FoSAM-Stream is outlined in Algorithm~\ref{alg:reuse}, a simple criterion is applied to detect if the view is still in the previous segment by calculating the difference between $I^{t}$ and $I{l}$ (line 4). If this difference exceeds a predefined threshold $\alpha$, it indicates a significant change in the front scene, triggering a full re-execution of XXXNet (line 5). In addition, if a saccade is detected (line 9), instance segmentation for the current frame can be skipped due to the reduced sensitivity of the human visual system during a saccade (line 10). If not, the current gaze location is analyzed to determine if it remains within the same IOI (line 13). If a predefined threshold $\beta$ is exceeded for the gaze location difference, XXXNet must be executed with the new gaze location, otherwise the last segmentation label map $Y^{last}$ can be reused.

\subsection{Impact of Adaptive Token Sampling Hyperparameters}
\textcolor{blue}{TODO}
\begin{wraptable}{r}{0.3\textwidth}
    \centering\resizebox{0.3\columnwidth}{!}{
       \begin{tabular}{c|c|c|c}
        \toprule
        $\tau$ & $c$ & FoSAM-S & Token Budget  \\ \midrule
        0.2 & 300 & 0.404 & 400 \\
        0.12 & 200 & 0.417 & 400\\
        0.1 & 100 & 0.416 & 400\\ \midrule
        0.4 & 150 & 0.39 & 200\\
        0.3 & 100 & 0.404 & 200\\
        0.25 & 50 & 0.402 & 200\\ \midrule
        0.7 & 75 & 0.355 & 100\\
        0.6 & 50 & 0.362 & 100\\
        0.55 & 25 & 0.362 & 100\\
        \bottomrule
        \end{tabular}
}
\caption{Evaluation results of adaptive token sampling settings.}
\label{tab:ablation:hyper}
\end{wraptable}

The saliency score threshold \(\tau\) and the minimum token count \(c\) are used to determine how many tokens are retained, as detailed in Section~\ref{sec:dynamic-window}. Different combinations of \(\tau\) and \(c\) can result in the same average number of tokens being retained across the samples in a dataset. In Table~\ref{tab:ablation:hyper}, we evaluate the effects of varying \(\tau\) and \(c\) under different token budgets using FoSAM-S over the Cityscape dataset. A larger \(\tau\) reduces the number of tokens selected based on saliency scores; to maintain the target token budget, \(c\) must be increased accordingly. However, a larger \(c\) restricts FoSAM’s ability to allocate fewer tokens for small IOIs. For example, under a token budget of 200, the configuration \(\tau = 0.3, c = 100\) achieves better performance than other settings.


\subsection{Impact of Loss Weights}
\textcolor{blue}{TODO}

During the training of FoSAM, the overall loss is computed as a weighted sum of the negative log-likelihood loss, classification loss, and Dice loss, with weights \(\lambda_1\), \(\lambda_2\), and \(\lambda_3\), respectively, as defined in Equation~\ref{eqn:loss1}, Equation~\ref{eqn:loss2}, and Equation~\ref{eqn:loss3}. Dice loss plays a critical role in the segmentation task, while the negative log-likelihood loss and classification loss help the saliency encoder in FoSAM better adapt the saliency score distribution based on the IOI. In Table~\ref{tab:ablation:loss}, we evaluate the contributions of the negative log-likelihood loss and the classification loss to the overall performance of FoSAM by varying their weights \(\lambda_1\) and \(\lambda_2\), using the FoSAM-S model on the Cityscapes dataset. 

\begin{wraptable}{r}{0.3\textwidth}
    \centering\resizebox{0.3\columnwidth}{!}{
       \begin{tabular}{c|c|c|c}
        \toprule
        $\lambda_1$ & $\lambda_2$ & Token Budget & IoU  \\ \midrule
        1 & 1 & 200 & 0.404 \\
        2 & 1 & 200 & 0.398 \\
        1 & 0 & 200 & 0.381 \\

        \bottomrule
        \end{tabular}
}
\caption{Impact of the weights of the nll loss ($\lambda_1$) and classification loss ($\lambda_2$).}
\label{tab:ablation:loss}
\end{wraptable}

The setting \(\lambda_1 = 1, \lambda_2 = 1\) is used as the baseline for our main evaluation results. Increasing \(\lambda_1\) reduces the relative contribution of the classification loss, which leads to a decrease in segmentation accuracy. When \(\lambda_2 = 0\), the classification loss is completely removed, resulting in a further drop in accuracy. This highlights the importance of the classification loss in achieving good segmentation performance.

\subsection{User Study}
\label{sec:userstudy_app}

\begin{wrapfigure}{r}{0.5\textwidth}
\vspace{-10pt}
\centering \includegraphics[width=0.5\textwidth]{figures/userstudy_sampleimage.pdf}
\caption{Four sample images.}
\label{fig:sample-image}
\vspace{-10pt}
\end{wrapfigure}

To evaluate the enhancement in user experience offered by FoSAM compared to the conventional SAM approach, we simulate their visual effects on the Meta Quest Pro HMD~\cite{quest_pro} due to HMD's restrictions on direct access to its NPU~\cite{qualcomm_hexagon_npu}. 

For each test image, segmentation masks for the IOIs were precomputed using FoSAM-S and SAM-B~\cite{kirillov2023segment}. During the user study, the HMD displays the mask of the currently observed IOI according to gaze locations obtained by the HMD's gaze tracker, simulating the visual effects of both algorithms.

To account for system performance, we profile the latency of FoSAM-S (200 token budget) and SAM-B on the Qualcomm XR2 Gen2 platform~\cite{qualcomm_xr2_gen2} using the Qualcomm AI Hub toolkit~\cite{qualcomm_aihub}, and artificially introduce the delay between the moment the user's gaze identifies the IOI and when the segmentation mask is displayed on the VR screen. Due to its efficient design, FoSAM-S exhibited significantly lower latency (21.6 ms) compared to SAM-B (1897.4 ms). Figure~\ref{fig:sample-image} illustrates these differences: in Image 1, FoSAM maintains low latency, ensuring masks closely align with current gaze location, whereas in Image 2, SAM’s higher latency results in noticeable misalignment between the mask and current gaze location, negatively impacting user experience.

Seven participants take part in the study (in Figure~\ref{fig:userstudy}, where the computer monitor display the HMD-cast content), and interact using the HMD controllers. The stimuli consists of four images (Figure~\ref{fig:sample-image}). The two methods, denoted as $m_1$ (FoSAM-S) and $m_2$ (SAM-B), are directly compared. Participants perform a two-interval forced-choice (2IFC) task~\cite{yeshurun2008bias}, viewing two segmentation results ($t_1$ and $t_2$) per image, with masks and latency applied as described. The test conditions ($t_1$ and $t_2$) are randomly assigned to $m_1$ and $m_2$. Using the HMD controller buttons, participants switch between $t_1$ and $t_2$ while observing different objects. After comparing both at least once, they select the one with higher perceived quality. Each participant completes 32 trials, consisting of 4 images, each tested with both $t$ and $m$ pairs across 4 repetitions, presented in a random sequence.

% \begin{figure}
%     \centering
%     \includegraphics[width=\linewidth]{figures/userstudy_sampleimages.pdf}
%     \caption{Four screenshots captured during the user study. In image 1 and 2, red dots indicate current gaze locations, blue cross represents the gaze location used for segmentation, reflecting a prior state due to latency.}
%     \label{fig:sample-image}
% \end{figure}

Figure~\ref{fig:userstudy-results} presents the results. Across participants, FoSAM was preferred in 96.9\%$\pm$4.8\% of trials overall, consistently outperforming SAM across all images (98.2\%$\pm$4.7\% for Image 1, 96.4\%$\pm$6.1\% for Image 2, 96.4\%$\pm$9.4\% for Image 3, and 96.4\%$\pm$9.4\% for Image 4). These results evidence that the superior performance of FoSAM can potentially improve user experience for AR/VR applications due to lower processing latency and precise segmentation accuracy.
\end{comment}
%%%%%%%%%%%%%%%%%%%%%%%%%%%%%%%%%%%%%%%%%%%%%%%%%%%%%%%%%%%%

\newpage
\section*{NeurIPS Paper Checklist}

%%% BEGIN INSTRUCTIONS %%%
The checklist is designed to encourage best practices for responsible machine learning research, addressing issues of reproducibility, transparency, research ethics, and societal impact. Do not remove the checklist: {\bf The papers not including the checklist will be desk rejected.} The checklist should follow the references and follow the (optional) supplemental material.  The checklist does NOT count towards the page
limit. 

Please read the checklist guidelines carefully for information on how to answer these questions. For each question in the checklist:
\begin{itemize}
    \item You should answer \answerYes{}, \answerNo{}, or \answerNA{}.
    \item \answerNA{} means either that the question is Not Applicable for that particular paper or the relevant information is Not Available.
    \item Please provide a short (1–2 sentence) justification right after your answer (even for NA). 
   % \item {\bf The papers not including the checklist will be desk rejected.}
\end{itemize}

{\bf The checklist answers are an integral part of your paper submission.} They are visible to the reviewers, area chairs, senior area chairs, and ethics reviewers. You will be asked to also include it (after eventual revisions) with the final version of your paper, and its final version will be published with the paper.

The reviewers of your paper will be asked to use the checklist as one of the factors in their evaluation. While "\answerYes{}" is generally preferable to "\answerNo{}", it is perfectly acceptable to answer "\answerNo{}" provided a proper justification is given (e.g., "error bars are not reported because it would be too computationally expensive" or "we were unable to find the license for the dataset we used"). In general, answering "\answerNo{}" or "\answerNA{}" is not grounds for rejection. While the questions are phrased in a binary way, we acknowledge that the true answer is often more nuanced, so please just use your best judgment and write a justification to elaborate. All supporting evidence can appear either in the main paper or the supplemental material, provided in appendix. If you answer \answerYes{} to a question, in the justification please point to the section(s) where related material for the question can be found.

IMPORTANT, please:
\begin{itemize}
    \item {\bf Delete this instruction block, but keep the section heading ``NeurIPS Paper Checklist"},
    \item  {\bf Keep the checklist subsection headings, questions/answers and guidelines below.}
    \item {\bf Do not modify the questions and only use the provided macros for your answers}.
\end{itemize} 
 

%%% END INSTRUCTIONS %%%


\begin{enumerate}

\item {\bf Claims}
    \item[] Question: Do the main claims made in the abstract and introduction accurately reflect the paper's contributions and scope?
    \item[] Answer: \answerTODO{} % Replace by \answerYes{}, \answerNo{}, or \answerNA{}.
    \item[] Justification: \justificationTODO{}
    \item[] Guidelines:
    \begin{itemize}
        \item The answer NA means that the abstract and introduction do not include the claims made in the paper.
        \item The abstract and/or introduction should clearly state the claims made, including the contributions made in the paper and important assumptions and limitations. A No or NA answer to this question will not be perceived well by the reviewers. 
        \item The claims made should match theoretical and experimental results, and reflect how much the results can be expected to generalize to other settings. 
        \item It is fine to include aspirational goals as motivation as long as it is clear that these goals are not attained by the paper. 
    \end{itemize}

\item {\bf Limitations}
    \item[] Question: Does the paper discuss the limitations of the work performed by the authors?
    \item[] Answer: \answerTODO{} % Replace by \answerYes{}, \answerNo{}, or \answerNA{}.
    \item[] Justification: \justificationTODO{}
    \item[] Guidelines:
    \begin{itemize}
        \item The answer NA means that the paper has no limitation while the answer No means that the paper has limitations, but those are not discussed in the paper. 
        \item The authors are encouraged to create a separate "Limitations" section in their paper.
        \item The paper should point out any strong assumptions and how robust the results are to violations of these assumptions (e.g., independence assumptions, noiseless settings, model well-specification, asymptotic approximations only holding locally). The authors should reflect on how these assumptions might be violated in practice and what the implications would be.
        \item The authors should reflect on the scope of the claims made, e.g., if the approach was only tested on a few datasets or with a few runs. In general, empirical results often depend on implicit assumptions, which should be articulated.
        \item The authors should reflect on the factors that influence the performance of the approach. For example, a facial recognition algorithm may perform poorly when image resolution is low or images are taken in low lighting. Or a speech-to-text system might not be used reliably to provide closed captions for online lectures because it fails to handle technical jargon.
        \item The authors should discuss the computational efficiency of the proposed algorithms and how they scale with dataset size.
        \item If applicable, the authors should discuss possible limitations of their approach to address problems of privacy and fairness.
        \item While the authors might fear that complete honesty about limitations might be used by reviewers as grounds for rejection, a worse outcome might be that reviewers discover limitations that aren't acknowledged in the paper. The authors should use their best judgment and recognize that individual actions in favor of transparency play an important role in developing norms that preserve the integrity of the community. Reviewers will be specifically instructed to not penalize honesty concerning limitations.
    \end{itemize}

\item {\bf Theory assumptions and proofs}
    \item[] Question: For each theoretical result, does the paper provide the full set of assumptions and a complete (and correct) proof?
    \item[] Answer: \answerTODO{} % Replace by \answerYes{}, \answerNo{}, or \answerNA{}.
    \item[] Justification: \justificationTODO{}
    \item[] Guidelines:
    \begin{itemize}
        \item The answer NA means that the paper does not include theoretical results. 
        \item All the theorems, formulas, and proofs in the paper should be numbered and cross-referenced.
        \item All assumptions should be clearly stated or referenced in the statement of any theorems.
        \item The proofs can either appear in the main paper or the supplemental material, but if they appear in the supplemental material, the authors are encouraged to provide a short proof sketch to provide intuition. 
        \item Inversely, any informal proof provided in the core of the paper should be complemented by formal proofs provided in appendix or supplemental material.
        \item Theorems and Lemmas that the proof relies upon should be properly referenced. 
    \end{itemize}

    \item {\bf Experimental result reproducibility}
    \item[] Question: Does the paper fully disclose all the information needed to reproduce the main experimental results of the paper to the extent that it affects the main claims and/or conclusions of the paper (regardless of whether the code and data are provided or not)?
    \item[] Answer: \answerTODO{} % Replace by \answerYes{}, \answerNo{}, or \answerNA{}.
    \item[] Justification: \justificationTODO{}
    \item[] Guidelines:
    \begin{itemize}
        \item The answer NA means that the paper does not include experiments.
        \item If the paper includes experiments, a No answer to this question will not be perceived well by the reviewers: Making the paper reproducible is important, regardless of whether the code and data are provided or not.
        \item If the contribution is a dataset and/or model, the authors should describe the steps taken to make their results reproducible or verifiable. 
        \item Depending on the contribution, reproducibility can be accomplished in various ways. For example, if the contribution is a novel architecture, describing the architecture fully might suffice, or if the contribution is a specific model and empirical evaluation, it may be necessary to either make it possible for others to replicate the model with the same dataset, or provide access to the model. In general. releasing code and data is often one good way to accomplish this, but reproducibility can also be provided via detailed instructions for how to replicate the results, access to a hosted model (e.g., in the case of a large language model), releasing of a model checkpoint, or other means that are appropriate to the research performed.
        \item While NeurIPS does not require releasing code, the conference does require all submissions to provide some reasonable avenue for reproducibility, which may depend on the nature of the contribution. For example
        \begin{enumerate}
            \item If the contribution is primarily a new algorithm, the paper should make it clear how to reproduce that algorithm.
            \item If the contribution is primarily a new model architecture, the paper should describe the architecture clearly and fully.
            \item If the contribution is a new model (e.g., a large language model), then there should either be a way to access this model for reproducing the results or a way to reproduce the model (e.g., with an open-source dataset or instructions for how to construct the dataset).
            \item We recognize that reproducibility may be tricky in some cases, in which case authors are welcome to describe the particular way they provide for reproducibility. In the case of closed-source models, it may be that access to the model is limited in some way (e.g., to registered users), but it should be possible for other researchers to have some path to reproducing or verifying the results.
        \end{enumerate}
    \end{itemize}


\item {\bf Open access to data and code}
    \item[] Question: Does the paper provide open access to the data and code, with sufficient instructions to faithfully reproduce the main experimental results, as described in supplemental material?
    \item[] Answer: \answerTODO{} % Replace by \answerYes{}, \answerNo{}, or \answerNA{}.
    \item[] Justification: \justificationTODO{}
    \item[] Guidelines:
    \begin{itemize}
        \item The answer NA means that paper does not include experiments requiring code.
        \item Please see the NeurIPS code and data submission guidelines (\url{https://nips.cc/public/guides/CodeSubmissionPolicy}) for more details.
        \item While we encourage the release of code and data, we understand that this might not be possible, so “No” is an acceptable answer. Papers cannot be rejected simply for not including code, unless this is central to the contribution (e.g., for a new open-source benchmark).
        \item The instructions should contain the exact command and environment needed to run to reproduce the results. See the NeurIPS code and data submission guidelines (\url{https://nips.cc/public/guides/CodeSubmissionPolicy}) for more details.
        \item The authors should provide instructions on data access and preparation, including how to access the raw data, preprocessed data, intermediate data, and generated data, etc.
        \item The authors should provide scripts to reproduce all experimental results for the new proposed method and baselines. If only a subset of experiments are reproducible, they should state which ones are omitted from the script and why.
        \item At submission time, to preserve anonymity, the authors should release anonymized versions (if applicable).
        \item Providing as much information as possible in supplemental material (appended to the paper) is recommended, but including URLs to data and code is permitted.
    \end{itemize}


\item {\bf Experimental setting/details}
    \item[] Question: Does the paper specify all the training and test details (e.g., data splits, hyperparameters, how they were chosen, type of optimizer, etc.) necessary to understand the results?
    \item[] Answer: \answerTODO{} % Replace by \answerYes{}, \answerNo{}, or \answerNA{}.
    \item[] Justification: \justificationTODO{}
    \item[] Guidelines:
    \begin{itemize}
        \item The answer NA means that the paper does not include experiments.
        \item The experimental setting should be presented in the core of the paper to a level of detail that is necessary to appreciate the results and make sense of them.
        \item The full details can be provided either with the code, in appendix, or as supplemental material.
    \end{itemize}

\item {\bf Experiment statistical significance}
    \item[] Question: Does the paper report error bars suitably and correctly defined or other appropriate information about the statistical significance of the experiments?
    \item[] Answer: \answerTODO{} % Replace by \answerYes{}, \answerNo{}, or \answerNA{}.
    \item[] Justification: \justificationTODO{}
    \item[] Guidelines:
    \begin{itemize}
        \item The answer NA means that the paper does not include experiments.
        \item The authors should answer "Yes" if the results are accompanied by error bars, confidence intervals, or statistical significance tests, at least for the experiments that support the main claims of the paper.
        \item The factors of variability that the error bars are capturing should be clearly stated (for example, train/test split, initialization, random drawing of some parameter, or overall run with given experimental conditions).
        \item The method for calculating the error bars should be explained (closed form formula, call to a library function, bootstrap, etc.)
        \item The assumptions made should be given (e.g., Normally distributed errors).
        \item It should be clear whether the error bar is the standard deviation or the standard error of the mean.
        \item It is OK to report 1-sigma error bars, but one should state it. The authors should preferably report a 2-sigma error bar than state that they have a 96\% CI, if the hypothesis of Normality of errors is not verified.
        \item For asymmetric distributions, the authors should be careful not to show in tables or figures symmetric error bars that would yield results that are out of range (e.g. negative error rates).
        \item If error bars are reported in tables or plots, The authors should explain in the text how they were calculated and reference the corresponding figures or tables in the text.
    \end{itemize}

\item {\bf Experiments compute resources}
    \item[] Question: For each experiment, does the paper provide sufficient information on the computer resources (type of compute workers, memory, time of execution) needed to reproduce the experiments?
    \item[] Answer: \answerTODO{} % Replace by \answerYes{}, \answerNo{}, or \answerNA{}.
    \item[] Justification: \justificationTODO{}
    \item[] Guidelines:
    \begin{itemize}
        \item The answer NA means that the paper does not include experiments.
        \item The paper should indicate the type of compute workers CPU or GPU, internal cluster, or cloud provider, including relevant memory and storage.
        \item The paper should provide the amount of compute required for each of the individual experimental runs as well as estimate the total compute. 
        \item The paper should disclose whether the full research project required more compute than the experiments reported in the paper (e.g., preliminary or failed experiments that didn't make it into the paper). 
    \end{itemize}
    
\item {\bf Code of ethics}
    \item[] Question: Does the research conducted in the paper conform, in every respect, with the NeurIPS Code of Ethics \url{https://neurips.cc/public/EthicsGuidelines}?
    \item[] Answer: \answerTODO{} % Replace by \answerYes{}, \answerNo{}, or \answerNA{}.
    \item[] Justification: \justificationTODO{}
    \item[] Guidelines:
    \begin{itemize}
        \item The answer NA means that the authors have not reviewed the NeurIPS Code of Ethics.
        \item If the authors answer No, they should explain the special circumstances that require a deviation from the Code of Ethics.
        \item The authors should make sure to preserve anonymity (e.g., if there is a special consideration due to laws or regulations in their jurisdiction).
    \end{itemize}


\item {\bf Broader impacts}
    \item[] Question: Does the paper discuss both potential positive societal impacts and negative societal impacts of the work performed?
    \item[] Answer: \answerTODO{} % Replace by \answerYes{}, \answerNo{}, or \answerNA{}.
    \item[] Justification: \justificationTODO{}
    \item[] Guidelines:
    \begin{itemize}
        \item The answer NA means that there is no societal impact of the work performed.
        \item If the authors answer NA or No, they should explain why their work has no societal impact or why the paper does not address societal impact.
        \item Examples of negative societal impacts include potential malicious or unintended uses (e.g., disinformation, generating fake profiles, surveillance), fairness considerations (e.g., deployment of technologies that could make decisions that unfairly impact specific groups), privacy considerations, and security considerations.
        \item The conference expects that many papers will be foundational research and not tied to particular applications, let alone deployments. However, if there is a direct path to any negative applications, the authors should point it out. For example, it is legitimate to point out that an improvement in the quality of generative models could be used to generate deepfakes for disinformation. On the other hand, it is not needed to point out that a generic algorithm for optimizing neural networks could enable people to train models that generate Deepfakes faster.
        \item The authors should consider possible harms that could arise when the technology is being used as intended and functioning correctly, harms that could arise when the technology is being used as intended but gives incorrect results, and harms following from (intentional or unintentional) misuse of the technology.
        \item If there are negative societal impacts, the authors could also discuss possible mitigation strategies (e.g., gated release of models, providing defenses in addition to attacks, mechanisms for monitoring misuse, mechanisms to monitor how a system learns from feedback over time, improving the efficiency and accessibility of ML).
    \end{itemize}
    
\item {\bf Safeguards}
    \item[] Question: Does the paper describe safeguards that have been put in place for responsible release of data or models that have a high risk for misuse (e.g., pretrained language models, image generators, or scraped datasets)?
    \item[] Answer: \answerTODO{} % Replace by \answerYes{}, \answerNo{}, or \answerNA{}.
    \item[] Justification: \justificationTODO{}
    \item[] Guidelines:
    \begin{itemize}
        \item The answer NA means that the paper poses no such risks.
        \item Released models that have a high risk for misuse or dual-use should be released with necessary safeguards to allow for controlled use of the model, for example by requiring that users adhere to usage guidelines or restrictions to access the model or implementing safety filters. 
        \item Datasets that have been scraped from the Internet could pose safety risks. The authors should describe how they avoided releasing unsafe images.
        \item We recognize that providing effective safeguards is challenging, and many papers do not require this, but we encourage authors to take this into account and make a best faith effort.
    \end{itemize}

\item {\bf Licenses for existing assets}
    \item[] Question: Are the creators or original owners of assets (e.g., code, data, models), used in the paper, properly credited and are the license and terms of use explicitly mentioned and properly respected?
    \item[] Answer: \answerTODO{} % Replace by \answerYes{}, \answerNo{}, or \answerNA{}.
    \item[] Justification: \justificationTODO{}
    \item[] Guidelines:
    \begin{itemize}
        \item The answer NA means that the paper does not use existing assets.
        \item The authors should cite the original paper that produced the code package or dataset.
        \item The authors should state which version of the asset is used and, if possible, include a URL.
        \item The name of the license (e.g., CC-BY 4.0) should be included for each asset.
        \item For scraped data from a particular source (e.g., website), the copyright and terms of service of that source should be provided.
        \item If assets are released, the license, copyright information, and terms of use in the package should be provided. For popular datasets, \url{paperswithcode.com/datasets} has curated licenses for some datasets. Their licensing guide can help determine the license of a dataset.
        \item For existing datasets that are re-packaged, both the original license and the license of the derived asset (if it has changed) should be provided.
        \item If this information is not available online, the authors are encouraged to reach out to the asset's creators.
    \end{itemize}

\item {\bf New assets}
    \item[] Question: Are new assets introduced in the paper well documented and is the documentation provided alongside the assets?
    \item[] Answer: \answerTODO{} % Replace by \answerYes{}, \answerNo{}, or \answerNA{}.
    \item[] Justification: \justificationTODO{}
    \item[] Guidelines:
    \begin{itemize}
        \item The answer NA means that the paper does not release new assets.
        \item Researchers should communicate the details of the dataset/code/model as part of their submissions via structured templates. This includes details about training, license, limitations, etc. 
        \item The paper should discuss whether and how consent was obtained from people whose asset is used.
        \item At submission time, remember to anonymize your assets (if applicable). You can either create an anonymized URL or include an anonymized zip file.
    \end{itemize}

\item {\bf Crowdsourcing and research with human subjects}
    \item[] Question: For crowdsourcing experiments and research with human subjects, does the paper include the full text of instructions given to participants and screenshots, if applicable, as well as details about compensation (if any)? 
    \item[] Answer: \answerTODO{} % Replace by \answerYes{}, \answerNo{}, or \answerNA{}.
    \item[] Justification: \justificationTODO{}
    \item[] Guidelines:
    \begin{itemize}
        \item The answer NA means that the paper does not involve crowdsourcing nor research with human subjects.
        \item Including this information in the supplemental material is fine, but if the main contribution of the paper involves human subjects, then as much detail as possible should be included in the main paper. 
        \item According to the NeurIPS Code of Ethics, workers involved in data collection, curation, or other labor should be paid at least the minimum wage in the country of the data collector. 
    \end{itemize}

\item {\bf Institutional review board (IRB) approvals or equivalent for research with human subjects}
    \item[] Question: Does the paper describe potential risks incurred by study participants, whether such risks were disclosed to the subjects, and whether Institutional Review Board (IRB) approvals (or an equivalent approval/review based on the requirements of your country or institution) were obtained?
    \item[] Answer: \answerTODO{} % Replace by \answerYes{}, \answerNo{}, or \answerNA{}.
    \item[] Justification: \justificationTODO{}
    \item[] Guidelines:
    \begin{itemize}
        \item The answer NA means that the paper does not involve crowdsourcing nor research with human subjects.
        \item Depending on the country in which research is conducted, IRB approval (or equivalent) may be required for any human subjects research. If you obtained IRB approval, you should clearly state this in the paper. 
        \item We recognize that the procedures for this may vary significantly between institutions and locations, and we expect authors to adhere to the NeurIPS Code of Ethics and the guidelines for their institution. 
        \item For initial submissions, do not include any information that would break anonymity (if applicable), such as the institution conducting the review.
    \end{itemize}

\item {\bf Declaration of LLM usage}
    \item[] Question: Does the paper describe the usage of LLMs if it is an important, original, or non-standard component of the core methods in this research? Note that if the LLM is used only for writing, editing, or formatting purposes and does not impact the core methodology, scientific rigorousness, or originality of the research, declaration is not required.
    %this research? 
    \item[] Answer: \answerTODO{} % Replace by \answerYes{}, \answerNo{}, or \answerNA{}.
    \item[] Justification: \justificationTODO{}
    \item[] Guidelines:
    \begin{itemize}
        \item The answer NA means that the core method development in this research does not involve LLMs as any important, original, or non-standard components.
        \item Please refer to our LLM policy (\url{https://neurips.cc/Conferences/2025/LLM}) for what should or should not be described.
    \end{itemize}

\end{enumerate}


\end{document}